% Untangling Federal Jurisdiction
% Part 1 of 3

% All content comes from Stephen Pratt. I have no idea if he approves of this effort or not.

\documentclass[xcolor=usenames,dvipsnames]{beamer}
\usetheme{Madrid}
%\usetheme{Goettingen}
%\usefonttheme{serif}
%\usefonttheme{structuresmallcapsserif}
% \usepackage[font=small,labelfont=bf]{caption}
\usefonttheme{professionalfonts}
\usepackage{xcolor}
\usepackage{rotating}
\usepackage{multirow}
\usepackage{soul}
\usepackage[overlay,absolute]{textpos}

\setbeamerfont{section title}{parent=title}
\setbeamercolor{section title}{parent=titlelike}
\defbeamertemplate*{my section page}{default}[1][]
{
    \centering
    \begin{beamercolorbox}[sep=8pt,center,#1]{section title}
        \usebeamerfont{section title}\insertsection\par
    \end{beamercolorbox}
}

\setbeamertemplate{navigation symbols}{}
\newcommand*{\mysectionpage}{\usebeamertemplate*{my section page}}

\def\Put(#1,#2)#3{\leavevmode\makebox(0,0){\put(#1,#2){#3}}}

\setbeamertemplate{footline}{}

\newenvironment<>{varblock}[2][\textwidth]{
    \begin{center}
        \begin{minipage}{#1}
            \setlength{\textwidth}{#1}
            \begin{actionenv}#3
                \def\insertblocktitle{#2}
                \par
                \usebeamertemplate{block begin}}
            {\par
                \usebeamertemplate{block end}
            \end{actionenv}
        \end{minipage}
    \end{center}
}

\setbeamertemplate{blocks}[default]
\addtobeamertemplate{block begin}{\pgfsetfillopacity{0.92}}{\pgfsetfillopacity{1}}
\newenvironment<>{quotepage}[3]{%
    \begin{frame}[t]%{#1}
        \centering
        %\begin{textblock}{100}[0.2,0](50,11)
            \includegraphics[height=\textheight,width=\textwidth,keepaspectratio]{#2} \\
        %\end{textblock}
        \begin{textblock}{70}[0,1](5,96)
            \begin{varblock}[0.9\paperwidth]{}%[shadow=false,rounded=true]{box2}
            %\colorbox{SeaGreen}{%
            %    \begin{minipage}[b]{0.95\paperwidth}
                    %\pgfsetfillopacity{0.35}
                    \begin{centering}
                    #3
                    % XXX Make this work correctly
                    \ifthenelse{\empty {#3} {}}{ \\ is there something here }{}
                    \end{centering}
                    \vspace{3pt}
                    \begin{minipage}[b]{\textwidth}
                    \hfill \small{#1}
                    \if{#1} \\ \fi
                    \end{minipage}
            %    \end{minipage}
            %}
            \end{varblock}
        \end{textblock}
        %\begin{textblock}
        % XXX Include credits (arg #1) here
        %\end{textblock}
}{%
    \end{frame}
}

\DeclareGraphicsExtensions{.pdf,.png,.jpg}

\setlength{\TPHorizModule}{0.01\paperwidth}
\setlength{\TPVertModule}{0.01\paperheight}
\textblockorigin{0mm}{0mm} % start everything near the top-left corner

\newcommand{\myquote}[3]{%
    \begin{quotepage}{#1}{#2}{#3}
    \end{quotepage}
}

%\usepackage{pgfpages} 
%\pgfpagesuselayout{resize to}[letterpaper,landscape] 

\usebackgroundtemplate%
{%
    \includegraphics[width=\paperwidth,height=\paperheight]{img/supcourt-facade-background.png}%
}

\begin{document}

%\title[Untangling Federal Jurisdiction]{Untangling Federal Jurisdiction Within a State}
%\author{Stephen Pratt}
%
%\frame{\titlepage}
%
%\section{A Recurrence to First Principles}
%
%\frame{\mysectionpage}

%\begin{frame}
%    \begin{columns}[onlytextwidth]
%        \column{0.5\textwidth}
%            \centering
%            \includegraphics{img/5000-year-leap.png}
%            \\ 28
%        \column{0.5\textwidth}
%            \pause
%            \centering
%            \includegraphics{img/making-of-america.png}
%            \\ 286
%    \end{columns}
%\end{frame}

\begin{frame}{Blackstone's Law Dictionary}
    \begin{columns}[onlytextwidth]
        \column{0.5\textwidth}
            \centering
            \includegraphics[width=0.75\textwidth]{img/blacks-law.png}
            \\ { \tiny 1920 pages }
        \column{0.5\textwidth}
        \textbf{Jurisdiction} \\
        Page 927 \\
        1.  A government's general power to exercise authority over all persons and things within its territory;
    \end{columns}
\end{frame}

\begin{frame}{Sheriff Gil Gilbertson}
    \begin{columns}[onlytextwidth]
        \column{0.5\textwidth}
            \centering
            Where does the United States Forest Service's authority come from?
            \\ { \tiny October, 2011 }
        \column{0.5\textwidth}
            \centering
            \includegraphics[width=0.75\textwidth]{img/gil-gilbertson.png}
            \\ Sheriff Gil Gilbertson
            \\ Josephine County, Oregon
    \end{columns}
\end{frame}

\begin{frame}{Sheriff Glenn Palmer}
    \begin{columns}[onlytextwidth]
        \column{0.5\textwidth}
            \centering
            \includegraphics[width=0.75\textwidth]{img/glenn-palmer-letter.png}
            \\ March 31, 2011
        \column{0.5\textwidth}
            \centering
            \includegraphics[width=0.75\textwidth]{img/glenn-palmer.png}
            \\ Sheriff Glenn Palmer
            \\ Grant County, Oregon
    \end{columns}
\end{frame}

\begin{frame}{Sheriff Glenn Palmer}
    \begin{columns}[onlytextwidth]
        \column{0.5\textwidth}
            \centering
            { \Large Where does the United States Forest Service get its
            \textbf{\emph{Constitutional}} authority to have law enforcement
            officers in Grant County? }
        \column{0.5\textwidth}
            \centering
            \includegraphics[width=0.75\textwidth]{img/glenn-palmer.png}
            \\ Sheriff Glenn Palmer
            \\ Grant County, Oregon
    \end{columns}
\end{frame}

\begin{frame}{Portland, Oregon}
    \includegraphics[width=0.3\textwidth]{img/doj.png}
    \includegraphics[width=0.3\textwidth]{img/usfs.png}
    \includegraphics[width=0.3\textwidth]{img/ossa.png}
    {   
        \centering
        \\ { \huge January 12, 2012 } \\
    }
\end{frame}

\begin{frame}{Portland, Oregon}
    \begin{columns}[onlytextwidth]
        \column{0.5\textwidth}
            \centering
            \includegraphics[height=0.28\textheight]{img/amanda-marshall.png}
            \\ United States Attorney \\ S. Amanda Marshall \\
            \includegraphics[height=0.28\textheight]{img/elmer-dickens.png}
            \\ Washington County Attorney \\ Elmer M. Dickens \\

        \column{0.5\textwidth}
            \centering
            \includegraphics[height=0.28\textheight]{img/gil-gilbertson.png}
            \\ Sheriff Gil Gilbertson \\ Josephine County \\
            \includegraphics[height=0.28\textheight]{img/glenn-palmer.png}
            \\ Sheriff Glenn Palmer \\ Grant County \\

    \end{columns}
\end{frame}

\begin{frame}
    \begin{columns}[onlytextwidth]
        \column{0.5\textwidth}
What bothered me the most regarding the conversation in Portland with the US
Attorneys office was that they did not want to talk about the Constitution,
but rather put all their confidence in the Supreme Court rulings. They even
referred to the ``Supreme Law of the Land'' as the Supreme Court and were visibly
upset that I would object to that notion
        \column{0.5\textwidth}
            \centering
            \includegraphics[width=0.75\textwidth]{img/gil-gilbertson.png}
            \\ Sheriff Gil Gilbertson
            \\ (Personal email, February 29, 2012)
    \end{columns}
\end{frame}

\def\braces#1{[#1]}

\begin{frame}{Summary and Opinion, January 13, 2012}
    \begin{columns}[onlytextwidth]
        \column{0.5\textwidth}
            \centering
            \includegraphics[width=0.75\textwidth]{img/elmer-dickens.png}
            \\ Attorney Elmer Dickens \\
        \column{0.5\textwidth}
\ldots it appears that the sheriff \braces{Gilbertson} does not believe that the US Supreme Court has the authority to interpret the meaning of the US Constitution\ldots
    \end{columns}
\end{frame}

\begin{frame}{Summary and Opinion, January 13, 2012}
    \begin{columns}[onlytextwidth]
        \column{0.5\textwidth}
        There are many Supreme Court cases going back over a hundred years\ldots \\
        \vspace{16pt}
        The federal Congress, well over one hundred years ago\ldots
        \column{0.5\textwidth}
            \centering
            \includegraphics[width=0.75\textwidth]{img/elmer-dickens.png}
            \\ Attorney Elmer Dickens \\
    \end{columns}
\end{frame}

\begin{frame}{Summary and Opinion, January 13, 2012}
    \begin{columns}[onlytextwidth]
        \column{0.5\textwidth}
To put it bluntly, the opinion that federal agents do not have jurisdiction on federal property located within the State of Oregon is without support in the law and is meritless.
        \column{0.5\textwidth}
            \centering
            \includegraphics[width=0.75\textwidth]{img/elmer-dickens.png}
            \\ Attorney Elmer Dickens \\
    \end{columns}
\end{frame}

\begin{frame}{Supreme Court, 2012}
    \centering
    {\LARGE How \emph{trustworthy} are the decisions of the Supreme Court? } \\
    \includegraphics[width=0.9\textwidth]{img/supreme.png} \\
\end{frame}

\begin{frame}{Justice Clarence Thomas}
    \begin{columns}[onlytextwidth]
        \column{0.5\textwidth}
            \centering
            \includegraphics[width=0.75\textwidth]{img/clarence-thomas.png} \\
        \column{0.5\textwidth}
        There are really only two ways to interpret the Constitution: try to
        discern as best we can what the framers intended or \emph{make it up}.
        Unless interpretive methodologies are tied to the
        \textbf{\emph{original intent}} of the framers, they have no more basis
        in the Constitution than the latest football scores. \\
        \vspace{10pt}
        Justice Clarence Thomas \\ Wall Street Journal Opinion \\ 20 October 2008
    \end{columns}
\end{frame}

\begin{frame}
    \centering
    \includegraphics[width=0.95\textwidth]{img/thomas-breyer.jpg} \\
    \pause
    { \large \ldots they just make it up! } \\
\end{frame}

\begin{frame}
    \begin{columns}[onlytextwidth]
        \column{0.5\textwidth}
            \centering
            \includegraphics[width=0.75\textwidth]{img/making-of-america.png} \\
            W. Cleon Skousen, Lawyer \\
            1985, 888 pages \\

        \column{0.5\textwidth}
            ``\ldots \braces{the Supreme Court} began reversing previous decisions by the bushel basket.''
    \end{columns}
\end{frame}

%\begin{frame}{Usurpation --- Timothy B. Lewis}
%    Show a decent shot of Timothy B. Lewis and/or his book \\
%    \includegraphics{img/timothy-lewis.png}
%\end{frame}

\begin{frame}{The Supremacists --- Phyllis Schlafly}
    \begin{columns}[onlytextwidth]
        \column{0.5\textwidth}
            \centering
            \includegraphics[width=0.75\textwidth]{img/the-supremacists.png} \\
            2004, 246 pages \\

        \column{0.5\textwidth}
            \centering
            \includegraphics[width=0.75\textwidth]{img/phyllis-schlafly.png} \\
            Lawyer \\
    \end{columns}
\end{frame}

\begin{frame}{Men in Black}
    \begin{columns}[onlytextwidth]
        \column{0.5\textwidth}
            \centering
            \includegraphics[width=0.75\textwidth]{img/men-in-black.png} \\

        \column{0.5\textwidth}
            \centering
            \includegraphics[width=0.75\textwidth]{img/mark-levin.png} \\
            Mark Levin, Lawyer \\
    \end{columns}
\end{frame}

\begin{frame}{U.S. Attorney S. Amanda Marshall}
    \begin{columns}[onlytextwidth]
        \column{0.5\textwidth}
            \centering
            \includegraphics[width=0.75\textwidth]{img/amanda-marshall.png} \\
            United States Attorney S. Amanda Marshall \\

        \column{0.5\textwidth}
            \centering
            \includegraphics[width=0.75\textwidth]{img/marshall-letter.png} \\
            January 23, 2012 \\
    \end{columns}
\end{frame}

\begin{frame}{U.S. Attorney S. Amanda Marshall}
    \begin{columns}[onlytextwidth]
        \column{0.5\textwidth}
            \centering
            \includegraphics[width=0.75\textwidth]{img/amanda-marshall.png} \\
            United States Attorney S. Amanda Marshall \\

        \column{0.5\textwidth}
            ``Supreme Court precedent for the past one hundred years is clear: under the Property Clause, Congress may delegate rulemaking authority to its federal agencies\ldots''
    \end{columns}
\end{frame}

\begin{frame}{The Supreme Court}
    \centering
    \includegraphics[width=0.75\textwidth]{img/supreme.png} \\
    How trustworty is Supreme Court precendent? How trustworthy is their interpretation of ``the property clause''?
\end{frame}

\begin{frame}{U.S. Attorney S. Amanda Marshall}
    \begin{columns}[onlytextwidth]
        \column{0.5\textwidth}
            \centering
            \includegraphics[width=0.75\textwidth]{img/amanda-marshall.png} \\
            United States Attorney S. Amanda Marshall \\

        \column{0.5\textwidth}
            \centering
            ``Your position is incorrect\ldots incorrect\ldots incorrect\ldots \\ erroneous\ldots \braces{You are on a} misguided crusade.'' \\
            \includegraphics[width=0.3\textwidth]{img/gil-gilbertson.png} \\
    \end{columns}
\end{frame}

\begin{frame}{U.S. Attorney S. Amanda Marshall}
    \begin{columns}[onlytextwidth]
        \column{0.5\textwidth}
            \centering
            \includegraphics[width=0.75\textwidth]{img/amanda-marshall.png} \\
            United States Attorney S. Amanda Marshall \\

        \column{0.5\textwidth}
            \centering
            ``\braces{It is} your constitutional duty to promote respect for federal laws in your county, particularly when validated by the Supreme Court, whether you agree with them or not.'' \\
            \vspace{10pt}
            \includegraphics[width=0.3\textwidth]{img/gil-gilbertson.png} \\
    \end{columns}
\end{frame}

\begin{frame}{The Oath?}
    \centering
    \includegraphics[width=0.7\textwidth]{img/oath.jpg} \\
    I swear to uphold and defend the federal laws and regulations, particularly when validated by the Supreme Court, whether I agree with them or not!
\end{frame}

\begin{frame}{The Oath?}
    \centering
    \includegraphics[width=0.7\textwidth]{img/oath-q.png} \\
    I swear to uphold and defend the federal laws and regulations, particularly when validated by the Supreme Court, whether I agree with them or not!
\end{frame}

\begin{frame}{Duty-bound}
    \centering
    \includegraphics[width=0.3\textwidth]{img/court-bldg.png}
    \includegraphics[width=0.3\textwidth]{img/cleveland.png}
            % XXX Andy Griffith
    \includegraphics[width=0.3\textwidth]{img/constitution.png} \\
    \vspace{18pt}
    By oath, is a sheriff duty-bound to question the constitutionality of federal laws and court decisions? \\
\end{frame}

\begin{frame}{Sheriff Cleveland}
    \begin{columns}[onlytextwidth]
        \column{0.5\textwidth}
            \centering
            \includegraphics[width=0.75\textwidth]{img/cleveland.png} \\

        \column{0.5\textwidth}
            \centering
            Sheriff Cleveland, 1871 \\
            Erie County, New York
    \end{columns}
\end{frame}

\begin{frame}{How Constitutional are ``Federal Laws''}
    \begin{columns}[onlytextwidth]
        \column{0.5\textwidth}
            \centering
            \includegraphics[width=0.75\textwidth]{img/cleveland2.png} \\

        \column{0.5\textwidth}
            Grover Cleveland was faithful to the Founder'’s Constitution and to
            what he believed were God'’s commandments and common sense.  He
            vetoed 584 Acts of Congress \\
            \pause
            \vspace{15pt}
            ``I can find no warrant for such an appropriation in the Constitution.'' \\

    \end{columns}
\end{frame}

\begin{frame}{How Constitutional are ``Federal Laws''}
    \begin{columns}[onlytextwidth]
        \column{0.5\textwidth}
            \centering
            \includegraphics[width=0.75\textwidth]{img/cleveland2.png} \\

        \column{0.5\textwidth}
            ``I ought to have a monument over me when I die, not for anything I have ever done, but for the foolishness I have put a stop to.''

    \end{columns}
\end{frame}

%\begin{frame}{How Constitutional are ``Federal Laws''}
%    \centering
%    { \Huge{The Freedom Index} } \\
%    \vspace{10pt}
%    \emph{A Congressional Scorecard Based on the U.S. Constitution} \\
%    \pause
%    January 9, 2012 \\
%    \vspace{15pt}
%    House / Senate Average Score on Ten Measures: \\
%    \vspace{20pt}
%    \pause
%    \color{red}{\Huge{44\%}} \\
%\end{frame}
%
%\begin{frame}{How Constitutional are ``Federal Laws''}
%    \centering
%    { \Large{United States Senate and House of Representatives }} \\
%    \vspace{20pt}
%    \includegraphics[width=0.75\textwidth]{img/house-senate-fail.png} \\
%\end{frame}
%
%\begin{frame}{Think Inside the Box}
%    \begin{columns}[onlytextwidth] \column{0.5\textwidth}
%            \centering
%            \includegraphics[width=0.75\textwidth]{img/box.png} \\
%
%        \column{0.5\textwidth}
%            \centering
%            \includegraphics[width=0.45\textwidth]{img/amanda-marshall.png} \\
%            \includegraphics[width=0.45\textwidth]{img/elmer-dickens.png} \\
%    \end{columns}
%\end{frame}
%
%\begin{frame}{Think Inside the Box}
%    \begin{columns}[onlytextwidth]
%        \column{0.5\textwidth}
%            \centering
%            \includegraphics[width=0.75\textwidth]{img/box.png} \\
%
%        \column{0.5\textwidth}
%            \begin{itemize}
%                \item Supreme Court decisions and acts of Congress ``over 100 years''
%                \pause
%                \item USC
%                \pause
%                \item Kleppe / FLPMA
%                \pause
%                \item Federal Register
%                \pause
%                \item Code of Federal Regulations
%            \end{itemize}
%    \end{columns}
%\end{frame}

\begin{frame}{Fundamental Principles}
    \begin{columns}[onlytextwidth]
        \column{0.5\textwidth}
            \centering
            \includegraphics[width=0.75\textwidth]{img/ben-franklin.png} \\
            Ben Franklin \\

        \column{0.5\textwidth}
            ``A frequent recurrence to \emph{fundamental principles}\ldots is absolutely necessary to preserve the blessings of liberty and keep a government free.''
    \end{columns}
\end{frame}

\begin{frame}{Fundamental Principles}
    \begin{columns}[onlytextwidth]
        \column{0.5\textwidth}
            \centering
            \includegraphics[width=0.75\textwidth]{img/patrick-henry.png} \\
            Patrick Henry \\

        \column{0.5\textwidth}
            ``No free government, or the blessing of liberty, can be preserved to any people but\ldots by a frequent recurrence to \emph{fundamental principles}.
    \end{columns}
\end{frame}

\begin{frame}{Fletcher vs. Peck, 10 U.S. 87}
    \centering
    \includegraphics[width=0.75\textwidth]{img/supreme-court-1810.png} \\
    Supreme Court, 1810 \\
\end{frame}

\begin{frame}{Fletcher vs. Peck, 10 U.S. 87}
    \begin{columns}[onlytextwidth]
        \column{0.5\textwidth}
            \centering
            \includegraphics[width=0.75\textwidth]{img/supreme-court-1810.png} \\
            Supreme Court, 1810 \\

        \column{0.5\textwidth}
            ``The security of a people against the misconduct of their rulers,
            must lie in the frequent recurrence to \emph{first principles}, and
            the imposition of adequate constitutional restrictions.''
    \end{columns}
\end{frame}

\begin{frame}
    \begin{columns}[onlytextwidth]
        \column{0.5\textwidth}
            ``The further back you can look, the farther forward you are likely to see.'' \\
            \vspace{20pt} Statehood, p. 19

        \column{0.5\textwidth}
            \centering
            \includegraphics[width=0.75\textwidth]{img/winston-churchill.png} \\
            Winston Churchill \\
    \end{columns}
\end{frame}

\begin{frame}{Roman Civil Law, 9 A. D.}
    \centering
    \includegraphics[width=0.75\textwidth]{img/europe-map.png} \\
    \only<1>{Roman Law \\}
%    \only<2>{Civil Law \\}
%    \only<3>{Roman Civil Law \\}
\end{frame}

\begin{frame}{Roman Law}
    \begin{columns}[onlytextwidth]
        \column{0.5\textwidth}
            \centering
            \includegraphics[height=0.55\textheight]{img/fasces.png} \\
            Fasces \\

        \column{0.5\textwidth}
            \begin{itemize}
                \item All powerful central government
                \item National sovereignty
                \item Tax and control everything
                \item Watchdog everyone's business
                \item Promise prosperity and greatness
                \item Fight wars anywhere on earth
                \pause
                \item \emph{Do whatever is necessary. No exceptions. No limits\ldots}
            \end{itemize}
    \end{columns}
\end{frame}

\begin{frame}[t]
    \centering
    \only<1>{\includegraphics[width=0.95\textwidth]{img/teutoburg.png} \\ }
    \only<2>{\includegraphics[width=0.95\textwidth]{img/Varus01.jpg} \\ }
    Battle of Teutoburg Forest, 9 A. D. \\
    Arminius of the Cherusci, Publius Quinctilius Varus \\
%    \only<2>{``We will rule our own territory!'' \\ }
\end{frame}

\begin{frame}{Germany, 1875 --- Hermanns Denkmal}
    \centering
    \includegraphics[width=0.75\textwidth]{img/hermans-denkmal.png} \\
\end{frame}

\begin{frame}{Germany, 1875 --- Hermanns Denkmal}
    \centering
    \includegraphics[width=0.75\textwidth]{img/herman2.png} \\
\end{frame}

%\begin{frame}
%    \begin{columns}[onlytextwidth]
%        \column{0.5\textwidth}
%            \centering
%            \includegraphics[width=0.75\textwidth]{img/herman2.png} \\
%
%        \column{0.5\textwidth}
%            Arminius / Herman preserved ``Common Law''
%            \pause
%            \begin{itemize}
%                \item Higher Law
%                \pause
%                \item Law of Nature
%                \pause
%                \item Constrast with Roman Law, based on coercion and force
%            \end{itemize}
%    \end{columns}
%\end{frame}
%
%\begin{frame}{Preserve Our Principles}
%    \centering
%    \includegraphics[width=0.75\textwidth]{img/herman3.png} \\
%    Rome, go home! \\
%\end{frame}

\begin{frame}{Anglo-Saxon Common Law}
    \begin{columns}[onlytextwidth]
        \column{0.5\textwidth}
            Germanic brothers who entered England in 450 A. D., bringing with them Anglo-Saxon Common Law, the ``Guardian of Freedom''.

        \column{0.5\textwidth}
            \centering
            \includegraphics[height=0.55\textheight]{img/hengist-horsa.png} \\
            Hengist and Horsa \\
    \end{columns}
\end{frame}

\begin{frame}{Fundamental Principles of Common Law}
    \centering
    \includegraphics[width=0.75\textwidth]{img/schoolroom.png} \\
\end{frame}

\begin{frame}{1215 A. D. --- Magna Carta}
    \centering
    Baronies force King John to sign the Magna Carta \\
    \includegraphics[width=0.75\textwidth]{img/king-john.png} \\
    ``\ldots to be sent to each of the English counties and to be read to all freemen.'' \\
\end{frame}

\begin{frame}{1215 A. D. --- Magna Carta}
    \begin{columns}[onlytextwidth]
        \column{0.5\textwidth}
            \centering
            \includegraphics[height=0.55\textheight]{img/magna-carta.png} \\

        \column{0.5\textwidth}
            ``All forests that have been created in our reign shall at once be disafforested.  River-banks that have been enclosed in our reign shall be treated similarly.'' \\
            (Clause 47)  \\
    \end{columns}
\end{frame}

\begin{frame}{1215 A. D. --- Magna Carta}
    \begin{columns}[onlytextwidth]
        \column{0.5\textwidth}
            \centering
            \includegraphics[height=0.55\textheight]{img/magna-carta.png} \\

        \column{0.5\textwidth}
            ``All evil customs relating to forest and warrens, foresters, warreners, sheriffs and their servants, or river-banks and their wardens, are at once to be investigated in every county  [and] are to be abolished completely and irrevocably.'' \\
            (Clause 48)
    \end{columns}
\end{frame}

%\begin{frame}{Preserve Our Principles}
%    \begin{columns}[onlytextwidth]
%        \column{0.5\textwidth}
%            \centering
%            \includegraphics[width=0.75\textwidth]{img/jefferson.png} \\
%            Thomas Jefferson \\
%
%        \column{0.5\textwidth}
%            Restore the Ancient Principles\ldots
%            \begin{itemize}
%                \item \ldots of Common Law
%                \item \ldots of the Old Testament
%            \end{itemize}
%    \end{columns}
%\end{frame}

\begin{frame}{Government of Ancient Israel}
    \centering
    \includegraphics[width=0.75\textwidth]{img/moses.jpg} \\
    \only<2>{ Groups of tens, fifties, hundreds, thousands, to successfully govern 3 million people \\ }
%    \only<3>{ This became the basis of ``Common'', or ``People's'' Law \\ }
\end{frame}

\begin{frame}{Anglo-Saxon Common Law}
    \begin{columns}[onlytextwidth]
        \column{0.5\textwidth}
            \centering
            \includegraphics[width=0.75\textwidth]{img/hengist-horsa.png} \\
%            Anglo-Saxon Common Law \color{red}\\

        \column{0.5\textwidth}
            \begin{itemize}
                \item Ten families --- a ``Tithing'', led by a ``tithing man''
                \item Two tithings --- a ``Vil'', led by a ``vil man''
                \item Two vils --- a ``Hundred'', led by a ``hundred man''
                \item Several hundreds --- a ``Shire'', led by a ``shire reef'', or later, ``sheriff''
            \end{itemize}
    \end{columns}
\end{frame}

\begin{frame}{Roman Law v. Common Law}
    \begin{columns}[onlytextwidth]
        \column{0.5\textwidth}
            \centering
            \includegraphics[height=0.75\textheight]{img/fasces.png} \\

        \column{0.5\textwidth}
            \centering
            \includegraphics[height=0.35\textheight]{img/hh-coin.png} \\
            \includegraphics[height=0.35\textheight]{img/israel-coin.png} \\
    \end{columns}
\end{frame}

\begin{frame}{Two Contending Forces}
    \begin{columns}[onlytextwidth]
        \column{0.5\textwidth}
            \begin{varblock}[0.93\textwidth]{}\centering \huge{Simple Republic} \\ \vspace{20pt} by coercion if necessary\end{varblock}

        \column{0.5\textwidth}
            \begin{varblock}[0.93\textwidth]{}\centering \huge{Compound Republic} \\ \vspace{20pt} voluntary membership\end{varblock}
    \end{columns}
    \textbf{\huge{ \color{red}
        \Put(155,120){v.}
    }}
\end{frame}

%\begin{frame}{``Shire Reef'' $\rightarrow$ ``Sheriff''}
%    \centering
%    \includegraphics[width=.9\textwidth]{img/shire-reef.png} \\
%\end{frame}
%
%\begin{frame}
%    \centering
%    \includegraphics[width=.9\textwidth]{img/sheriffs.png} \\
%    \pause
%    \Put(0.3,6){\textbf{\color{blue}{\Large{Maryland, 1634}}}}
%    \pause
%    \Put(3,0.8){\textbf{\color{blue}{\Large{Virginia, 1651}}}}
%\end{frame}

%\begin{frame}{Utah Territorial Act, September 9, 1850}
%    \centering
%    \begin{block}{}
%        ``\ldots and the same is hereby, created into a \emph{temporary government}, by the name of the Territory of Utah\ldots''
%    \end{block}
%    \vspace{10pt}
%    \includegraphics[width=0.75\textwidth]{img/utah-terr.png} \\
%\end{frame}
%
%\begin{frame}{Utah Territory, 1894}
%    \begin{columns}[onlytextwidth]
%        \column{0.5\textwidth}
%            \centering
%            \includegraphics[height=0.75\textheight]{img/utah-state.png} \\
%
%        \column{0.5\textwidth}
%            \Huge{People \\ + Territory \\ + Governance \\ = State}
%    \end{columns}
%\end{frame}
%
%\begin{frame}{Utah Territory, 1894}
%    \begin{columns}[onlytextwidth]
%        \column{0.5\textwidth}
%            \centering
%            \includegraphics[height=0.75\textheight]{img/utah-state.png} \\
%
%        \column{0.5\textwidth}
%            \begin{block}{Utah Enabling Act \textbf{Compact}, July 16, 1894}
%                As of May 3, 2012, it appears there are still \textbf{\emph{{\color{red}seven violations}}} of the Compact on the part of the Federal Government.
%            \end{block}
%    \end{columns}
%\end{frame}
%
%\begin{frame}{Utah Territory}
%    \begin{columns}[onlytextwidth]
%        \column{0.5\textwidth}
%            \centering
%            \includegraphics[height=0.75\textheight]{img/constitution.png} \\
%
%        \column{0.5\textwidth}
%            \begin{block}{Admissions Clause}
%                ``New States may be admitted by Congress into this Union\ldots''
%            \end{block}
%    \end{columns}
%\end{frame}
%
%\begin{frame}{Utah, January 4, 1896}
%    \centering
%    \includegraphics[height=0.75\textheight]{img/utah-constitution.png} \\
%\end{frame}

\begin{frame}{Jurisdiction Clause Family}
    \centering
    \includegraphics[height=0.85\textheight]{img/family.png} \\
    \huge{\textbf{ \color{white}
        \pause
        \Put(25, 300){\begin{sideways}\colorbox{blue}{Admissions}\end{sideways}}
        \pause
        \Put(60, 300){\begin{sideways}\colorbox{blue}{Claims}\end{sideways}}
        \pause
        \Put(90, 300){\begin{sideways}\colorbox{blue}{Property}\end{sideways}}
        \pause
        \Put(134, 300){\begin{sideways}\colorbox{blue}{Guarantee}\end{sideways}}
        \pause
        \Put(168, 300){\begin{sideways}\colorbox{blue}{Enclave}\end{sideways}}
        \pause
        \Put(200, 300){\begin{sideways}\colorbox{blue}{Engagements}\end{sideways}}
        \pause
        \Put(235, 300){\begin{sideways}\colorbox{blue}{Supremacy}\end{sideways}}
    }}
\end{frame}

%\unitlength=1in

%\begin{frame}{State of Utah}
%    \begin{columns}[onlytextwidth]
%        \column{0.5\textwidth}
%            \centering
%            \includegraphics[height=0.75\textheight]{img/utah-constitution.png} \\
%
%        \column{0.5\textwidth}
%            The new state legislature began passing laws to protect life, liberty, and property. \\
%            \vspace{20pt}
%            County Sheriff established by:
%            \begin{itemize}
%                \item 17-16-1
%                \item 17-22-1.5
%                \item 20A-9-201
%                \item 53-6-202(4)
%                \item 53-6-205
%            \end{itemize}
%    \end{columns}
%\end{frame}
%
%\begin{frame}
%    \begin{varblock}[0.8\textwidth]{Definition}
%        A legal dictionary defines the Sheriff as ``\ldots a creature of law created by the \textbf{sovereign power in the state}\ldots'' 
%    \end{varblock}
%\end{frame}

\begin{frame}{Mack / Printz v. The United States of America, 1997}
    \centering
    \includegraphics[width=0.65\textwidth]{img/printz-mack.jpg} \\
\end{frame}

\begin{frame}{Mack / Printz v. The United States of America, 1997}
    \begin{columns}[onlytextwidth]
        \column{0.5\textwidth}
            \centering
            \includegraphics[height=0.75\textheight]{img/scalia.jpg} \\
            Justice Antonin Scalia \\

        \column{0.5\textwidth}
            %\begin{block}{Mack/Printz v. USA}
            %    ``\ldots it is incontestible that the Constitution established a system of `dual sovereignty'.''
            %\end{block}
            \includegraphics[width=\textwidth]{img/incontestable.png} \\
            \pause
            {
                \color{red}
                \emph{Did the United States Constitution establish a system of ``dual sovereignty''?}
            }
    \end{columns}
\end{frame}

\begin{frame}{Mack / Printz v. The United States of America, 1997}
    \begin{columns}[onlytextwidth]
        \column{0.5\textwidth}
            \centering
            \includegraphics[height=0.75\textheight]{img/scalia.jpg} \\
            Justice Antonin Scalia \\

        \column{0.5\textwidth}
            %\begin{block}{Mack/Printz v. USA}
            %    ``Although the States surrendered many of their powers to the new Federal Government\ldots''
            %\end{block}
            \includegraphics[width=\textwidth]{img/surrendered.png} \\
            \pause
            {
                \color{red}
                \emph{Did the States ``[surrender] many of their powers to the new Federal Government''?}
            }
    \end{columns}
\end{frame}

\begin{frame}{Sovereignty}
\begin{table}[h]
\centering
\begin{tabular}{p{.48\textwidth} p{.48\textwidth}} 
    { \centering \huge{State's rights} \\ } & { \centering \huge{Consolidationists} \\ } \\
    { \centering \includegraphics[height=0.45\textheight]{img/noah-webster.jpg} \\ } &
        { \centering \includegraphics[height=0.45\textheight]{img/daniel-webster-portrait.png} \\ } \\
    \only<1>{{\centering Noah Webster \\ } & {\centering Daniel Webster \\ } \\ }
    \only<2>{ { \centering \textbf{Sovereignty:} Supreme in power; \ldots the highest power \\ } &
      { \centering \textbf{Sovereignty:} The sum of all rights and powers \\ } \\ }
\end{tabular}
\end{table}
\end{frame}

\begin{frame}{Summary}
    Sovereignty is possession of the highest power. It is impossible for both
    the states and the federal government to possess the highest power at the
    same time. When the sovereign states created the Constitution, they
    delegated a few enumerated powers to their new agent called the Federal
    Government. \\
    \vspace{20pt}
    \pause
    { \color{red}\textbf{\emph}They did not ``surrender many of their powers to the Federal Government''}
\end{frame}

\begin{frame}{``Delegate'' vs. ``Surrender''}
    \begin{itemize}
        \item Delegate = send with a trust or commission to act for another. Temporary.
        \item Surrender = to yield, to cede, to give up, to grant. Permanent.
    \end{itemize}
\end{frame}


\begin{frame}{Tea Party}
    \centering
    \includegraphics[width=0.75\textwidth]{img/tea-party.jpg} \\
    \begin{block}{}Current choices in Washington D. C. are ``a grave threat to our national sovereignty.'' \end{block}
    \begin{itemize}
        \item What is national sovereignty?
        \item Did the United States Constitution establish national sovereignty?
        \item What is state sovereignty?
    \end{itemize}
\end{frame}

\begin{frame}{Alexander Hamilton, Constitutional Convention}
    \begin{columns}[onlytextwidth]
        \column{0.5\textwidth}
            \centering
            \includegraphics[height=0.75\textheight]{img/hamilton.png} \\
            Alexander Hamilton \\

        \column{0.5\textwidth}
            \begin{block}{}
                ``Two sovereignties cannot co-exist within the same limits.''
            \end{block}
            \pause
            \vspace{15pt}
            In other words, \emph{state sovereignty} and \emph{national sovereignty} cannot exist at the same time.
    \end{columns}
\end{frame}

\begin{frame}{Alexander Hamilton, Constitutional Convention}
    \begin{columns}[onlytextwidth]
        \column{0.5\textwidth}
            \centering
            \includegraphics[height=0.75\textheight]{img/hamilton.png} \\
            Alexander Hamilton \\

        \column{0.5\textwidth}
            \begin{block}{}
                ``[We must establish] a complete sovereignty in the General Government\ldots The general power, if it preserves itself, must swallow up the state powers.''
            \end{block}
            \pause
            {
                \centering
                \includegraphics[width=0.5\textwidth]{img/red-light.png} \\
            }
    \end{columns}
\end{frame}

\begin{frame}{Thomas Jefferson, 1824}
    \begin{columns}[onlytextwidth]
        \column{0.5\textwidth}
            \centering
            \includegraphics[height=0.75\textheight]{img/jefferson.png} \\

        \column{0.5\textwidth}
            Men by their constitutions are naturally divided into two parties: \\
            \begin{enumerate}
                \item Those who fear and distrust the people, and wish to draw all powers from them into the hands of the higher classes.
                \item Those who identify themselves with the people, [and] have confidence in them.
            \end{enumerate}
    \end{columns}
\end{frame}

\begin{frame}{Thomas Jefferson, 1824}
    \begin{columns}[onlytextwidth]
        \column{0.5\textwidth}
            \centering
            \includegraphics[height=0.75\textheight]{img/jefferson.png} \\

        \column{0.5\textwidth}
            ``In every country these two parties exist.''
    \end{columns}
\end{frame}

%\begin{frame}{Two Contending Forces}
%    \begin{columns}[onlytextwidth]
%        \column{0.5\textwidth}
%            \begin{varblock}[0.9\textwidth]{}\huge{We the people of the United States is a nationalist system}\end{varblock}
%
%        \column{0.5\textwidth}
%            \begin{varblock}[0.9\textwidth]{}\huge{We the people of the States of NH, MA, RI, NY, PA, NJ, \ldots  are a federative system.}\end{varblock}
%    \end{columns}
%    \textbf{\huge{ \color{red}
%        \Put(155,150){v.}
%    }}
%\end{frame}
%
%
%\begin{frame}{Remember\ldots}
%    \begin{columns}[onlytextwidth]
%        \column{0.5\textwidth}
%            \centering
%            \huge{ROMAN LAW} \\
%            \vspace{20pt}
%            \includegraphics[height=0.55\textheight]{img/fasces.png} \\
%
%        \column{0.5\textwidth}
%            \begin{itemize}
%                \item All powerful central government
%                \item National sovereignty
%                \item Tax and control everything
%                \item Watchdog everyone's business
%                \item Promise prosperity and greatness
%                \item Fight wars anywhere on earth
%            \end{itemize}
%    \end{columns}
%\end{frame}
%
%\begin{frame}
%    \begin{columns}[onlytextwidth]
%        \column{0.5\textwidth}
%            \centering
%            \includegraphics[width=0.75\textwidth]{img/andrew-johnson.jpg} \\
%            Andrew Johnson \\
%
%        \column{0.5\textwidth}
%            \begin{block}{News Release, August 14, 1866}
%                The National Union Party declares ``the absolute supremacy of the National Government''
%            \end{block}
%    \end{columns}
%\end{frame}
%
%\begin{frame}{1914}
%    \centering
%    \includegraphics[width=.9\textwidth]{img/lincoln-memorial.jpg} \\
%    \pause
%    \Put(250,210){\includegraphics[height=0.55\textwidth]{img/fasces.png}}
%\end{frame}
%
%\begin{frame}
%    \begin{columns}[onlytextwidth]
%        \column{0.5\textwidth}
%            \centering
%            \includegraphics[width=0.75\textwidth]{img/americas-caesar.png} \\
%
%        \column{0.5\textwidth}
%            \begin{block}{America's Caesar}
%                The Decline and Fall of Republican Government in the United States of America
%            \end{block}
%    \end{columns}
%\end{frame}
%
%\begin{frame}
%    \begin{columns}[onlytextwidth]
%        \column{0.5\textwidth}
%            \centering
%            \includegraphics[width=0.75\textwidth]{img/dime.png} \\
%            1916 - 1945 \\
%
%        \column{0.5\textwidth}
%            \centering
%            \includegraphics[height=0.55\textheight]{img/fasces.png} \\
%            Roman Law\\
%
%    \end{columns}
%\end{frame}
%
%\begin{frame}
%    \centering
%    \includegraphics[width=.9\textwidth]{img/obama-fasces.jpg} \\
%    \pause
%    \Put(250,210){\includegraphics[height=0.55\textwidth]{img/fasces.png}}
%\end{frame}
%
%\begin{frame}
%    \centering
%    \includegraphics[width=.9\textwidth]{img/reich-stamp.png} \\
%\end{frame}
%
%\begin{frame}{Two Contending Forces}
%    \begin{columns}[onlytextwidth]
%        \column{0.5\textwidth}
%            \begin{varblock}[0.9\textwidth]{}\huge{Sovereignty of the National Government}\end{varblock}
%
%        \column{0.5\textwidth}
%            \begin{varblock}[0.9\textwidth]{}\huge{Sovereignty of the individual States}\end{varblock}
%    \end{columns}
%    \textbf{\huge{ \color{red}
%        \Put(155,120){v.}
%    }}
%\end{frame}
%
%

\begin{frame}
    \begin{columns}[onlytextwidth]
        \column{0.5\textwidth}
            \centering
            \includegraphics[width=0.75\textwidth]{img/upshur.jpg} \\
            Abel Parker Upshur \\

        \column{0.5\textwidth}
            ``In order to a correct understanding of the Constitution, it is absolutely necessary to understand the situation of the States before it was adopted.''
    \end{columns}
\end{frame}

\begin{frame}
    \begin{columns}[onlytextwidth]
        \column{0.5\textwidth}
            \centering
            \includegraphics[height=0.75\textheight]{img/colonies.png} \\

        \column{0.5\textwidth}
            \begin{block}{The situtation, in 1776}
            \textbf{Thirteen} separate, free, and independent \textbf{republics called ``states.''}
            \end{block}
    \end{columns}
\end{frame}

\begin{frame}{Sovereign States, 1788}
    \includegraphics[height=0.75\textheight]{img/colonies.png} \\
    \pause
    \Put(200,300){\includegraphics[width=0.4\textwidth]{img/constitution.png}}
    \Put(195,420){\textbf{List of delegated powers}}
    \pause
    \Put(200,60){\includegraphics[width=0.3\textwidth]{img/national-seal.png}}
    \Put(225,170){\textbf{\Large{Agent}}}
\end{frame}

\begin{frame}{James Madison, 1821}
    \begin{columns}[c]
        \column{0.5\textwidth}
            \centering
            \includegraphics[width=0.75\textwidth]{img/madison.jpg} \\
        \column{0.5\textwidth}
            ``Our government system is established by \emph{compact}, not between the Government of the United States and the State Governments, but between the \textbf{states as sovereign communities}.''
    \end{columns}
\end{frame}

\begin{frame}
    \begin{columns}[onlytextwidth]
        \column{0.5\textwidth}
            \centering
            \includegraphics[height=0.75\textheight]{img/walter-neal-book.jpg} \\
            1910,  p. 81\\

        \column{0.5\textwidth}
            \begin{itemize}
                \item In wresting the power, or what is called the sovereignty, from the crown, it passed directly to the people.
                \item It passed directly to the people of each colony, and not to the people of all the colonies in the aggregate --- to thirteen distinct communities, and not to one.
            \end{itemize}
    \end{columns}
\end{frame}

\begin{frame}{Second Continental Congress}
    \begin{columns}[onlytextwidth]
        \column{0.5\textwidth}
            The Articles of Confederation were first drafted by the Continental Congress in 1777 as ``a coordinating mechanism'', and sent to the states for ratification. They were finally ratified by each state in 1781.

        \column{0.5\textwidth}
            \centering
            \includegraphics[width=0.95\textwidth]{img/continental-congress.jpg} \\
    \end{columns}
\end{frame}

\begin{frame}{Second Continental Congress}
    \begin{columns}[onlytextwidth]
        \column{0.5\textwidth}
            \centering
            \includegraphics[width=0.95\textwidth]{img/continental-congress.jpg} \\

        \column{0.5\textwidth}
            \begin{block}{October 10, 1780}
                The Continental Congress passed the ``Resolution of 1780'' removing the only obstacle to a final ratification of the Articles of Confederation.
            \end{block}
    \end{columns}
\end{frame}

\begin{frame}
    \begin{alertblock}{Public lands}
        Beginning with The Resolution of 1780, the federal trust respecting public lands obligated the United States to extinguish both their governmental jurisdiction and their title to land that achieved statehood.
    \end{alertblock}
\end{frame}

\begin{frame}{Articles of Confederation}
    \begin{columns}[onlytextwidth]
        \column{0.5\textwidth}
            \centering
            \includegraphics[width=0.95\textwidth]{img/articles-of-confederation.png} \\

        \column{0.5\textwidth}
            \begin{block}{Article I, Section II}
                ``Each State retains its sovereignty, freedom and independence, and every power jurisdiction and right, which is not by this confederation expressly delegated to the United States, in Congress Assembled.''
            \end{block}
    \end{columns}
\end{frame}

\begin{frame}{Constitutional Convention of 1787}
    \centering
    \includegraphics[width=0.95\textwidth]{img/constitutional-convention.jpg} \\
\end{frame}

\begin{frame}{Constitutional Convention of 1787}
    \begin{columns}[onlytextwidth]
        \column{0.5\textwidth}
            \centering
            \includegraphics[width=0.95\textwidth]{img/constitutional-convention.jpg} \\

        \column{0.5\textwidth}
            Assignment:
            \begin{enumerate}
                \item Amend the Articles of Confederation
                \item Create a “national government”
            \end{enumerate}
    \end{columns}
\end{frame}

\begin{frame}{Constitutional Principles of Jurisdiction}
    \begin{columns}[onlytextwidth]
        \column{0.5\textwidth}
            \centering
            \includegraphics[width=0.95\textwidth]{img/constitution.png} \\

        \column{0.5\textwidth}
            \begin{itemize}
                \item ALL FEDERAL LAWS shall be made by the Congress of the United States
                \item The new Federal Government is an agent serving the states.
                \item The delegated powers are “few and defined”. 
                \item All powers not listed are retained by the states or the people.
                \item The “Resolution of 1780”  formed the basis upon which Congress was required to dispose of territorial and public lands.
            \end{itemize}
    \end{columns}
\end{frame}

%\section{Clause Family Martyred}
%\frame{\mysectionpage}

\begin{frame}{South Carolina v. United States, 1905}
    \centering
    \includegraphics[width=0.75\textwidth]{img/sc-1905.png} \\
    ``The Constitution is a written instrument. As such, its meaning does not alter. That which it meant when it was adopted, it means now.'' \\
\end{frame}

\begin{frame}{United States v. Sprague, 1931}
    \centering
    \includegraphics[width=0.75\textwidth]{img/sc-1905.png} \\
    ``The Constitution was written to be understood by the voters; its words and their phrases were used in their normal and ordinary\ldots meaning.'' \\
\end{frame}

\begin{frame}
    \begin{columns}[c]
        \column{0.5\textwidth}
            \centering
            \includegraphics[width=0.75\textwidth]{img/dictionary.png} \\
        \column{0.5\textwidth}
            \centering
            \includegraphics[height=0.5\textheight]{img/noah-webster.jpg} \\
            \only<1>{Published by Noah Webster in 1828 \\}
            \only<2>{\textbf{Engagement:} Obligation by agreement or contract. \\}
    \end{columns}
\end{frame}

\begin{frame}{Jurisdiction Clause Family --- Engagement}
    \centering
    \includegraphics[height=0.85\textheight]{img/family.png} \\
    \huge{\textbf{ \color{white}
        %\Put(25, 300){\begin{sideways}\colorbox{blue}{Admissions}\end{sideways}}
        %\Put(60, 300){\begin{sideways}\colorbox{blue}{Claims}\end{sideways}}
        %\Put(90, 300){\begin{sideways}\colorbox{blue}{Property}\end{sideways}}
        %\Put(134, 300){\begin{sideways}\colorbox{blue}{Guarantee}\end{sideways}}
        %\Put(168, 300){\begin{sideways}\colorbox{blue}{Enclave}\end{sideways}}
        \Put(200, 300){\begin{sideways}\colorbox{blue}{Engagements}\end{sideways}}
        %\Put(235, 300){\begin{sideways}\colorbox{blue}{Supremacy}\end{sideways}}
    }}
\end{frame}

\begin{frame}{Engagement Clause}
   \centering
   \includegraphics[height=.7\textheight]{img/constitution.png} \\
   \textbf{Article VI, clause 1:} ``\emph{Engagements} entered into, \emph{before} the Adoption of this Constitution, shall be as valid against the United States under this Constitution, as under the Confederation.''
\end{frame}

\begin{frame}{Articles of Confederation, on Federal Property}
    \begin{columns}[c]
        \column{0.5\textwidth}
            \centering
            \includegraphics[width=0.75\textwidth]{img/articles-of-confederation.png} \\
        \column{0.5\textwidth}
            \textbf{Article IX:} \\
            ``\ldots no State shall be deprived of territory for the benefit of the United States.''
    \end{columns}
\end{frame}

\begin{frame}
    \begin{block}{}
        Formation of a ``more perfect union'' does not absolve that union of prior engagements, including those obligations established by the Resolution of 1780 and the Articles of Confederation.''
    \end{block}
\end{frame}

\begin{frame}{Johnson v. McIntosh, 1823}
    \centering
    \includegraphics[width=0.75\textwidth]{img/sc-1905.png} \\
    ``And the \emph{vacant soil is to be disposed of} by that organ of the government which has the constitutional power to dispose of the national domain.'' \\
\end{frame}

\begin{frame}{Friedman v. Goodwin, 1856}
    \centering
    \includegraphics[width=0.75\textwidth]{img/sc-1905.png} \\
    ``On [California's] admission into the Union, \emph{she became the owner of all the public land} not disposed of by law of congress.'' \\
\end{frame}

%\begin{frame}
%    \begin{varblock}[0.8\textwidth]{Truism}
%        \Large{We may only delegate to government that which we have the right to do ourselves.}
%    \end{varblock}
%\end{frame}
%
%\begin{frame}{Two Contending Forces}
%    \begin{columns}[c]
%        \column{0.5\textwidth}
%            \centering
%            \large{Hamiltonian} \\
%            \vspace{15pt}
%            \includegraphics[height=0.60\textheight]{img/hamilton-portrait.png} \\
%        \column{0.5\textheight}
%            \centering
%            \large{Jeffersonian} \\
%            \vspace{15pt}
%            \includegraphics[height=0.60\textheight]{img/jefferson.png} \\
%    \end{columns}
%    \pause
%    \vspace{20pt}
%    On the 4th of July, we praise Jefferson, but we live in Hamilton's world.
%\end{frame}
%
%\begin{frame}
%    \begin{columns}[c]
%        \column{0.5\textwidth}
%            \centering
%            \includegraphics[height=0.60\textheight]{img/hamilton-portrait.png} \\
%            Alexander Hamilton \\
%        \column{0.5\textheight}
%            ``The greatest man who ever lived was Julius Caesar."
%    \end{columns}
%\end{frame}
%
%\begin{frame}{Two Contending Forces}
%    \begin{columns}[c]
%        \column{0.5\textwidth}
%            \centering
%            \large{Hamiltonian} \\
%            \vspace{15pt}
%            \includegraphics[height=0.60\textheight]{img/fasces.png} \\
%        \column{0.5\textheight}
%            \centering
%            \large{Jeffersonian} \\
%            \vspace{15pt}
%            \pause
%            \includegraphics[height=0.60\textheight]{img/hh-coin.png} \\
%    \end{columns}
%\end{frame}
%
%\begin{frame}{George Bancroft}
%    \begin{columns}[c]
%        \column{0.5\textwidth}
%            \centering
%            \includegraphics[width=0.75\textwidth]{img/george-bancroft.png} \\
%            George Bancroft, 1800 - 1891 \\
%        \column{0.5\textheight}
%            \centering
%            \textbf{A PLEA FOR THE CONSTITUTION OF THE UNITED STATES} \\
%            Wounded in the house of its guardians \\
%            \vspace{30pt} 1886, 25 pages \\
%    \end{columns}
%\end{frame}
%
%\begin{frame}{George Bancroft}
%    \begin{columns}[c]
%        \column{0.5\textwidth}
%            \centering
%            \includegraphics[width=0.75\textwidth]{img/george-bancroft.png} \\
%            George Bancroft, 1800 - 1891 \\
%        \column{0.5\textheight}
%            ``The government of the United States is far, very far removed from the powers of a sovereign state.''
%    \end{columns}
%\end{frame}
%
%\begin{frame}{``Consolidating School''}
%    \begin{columns}[c]
%        \column{0.5\textwidth}
%            \centering
%            \includegraphics[height=0.4\textheight]{img/story-portrait.png} \\
%            \includegraphics[height=0.4\textheight]{img/daniel-webster-portrait.png} \\
%        \column{0.5\textheight}
%            Sovereignty is the sum of all rights and powers. \\
%            In the Constitution sovereignty is delegated to the central government, which is an irrevocable \emph{ceding} of power.
%    \end{columns}
%\end{frame}
%
%\begin{frame}{``Consolidating School''}
%    \begin{columns}[c]
%        \column{0.5\textwidth}
%            \centering
%            \includegraphics[height=0.4\textheight]{img/story-portrait.png} \\
%            \includegraphics[height=0.4\textheight]{img/daniel-webster-portrait.png} \\
%        \column{0.5\textheight}
%            States have sovereignty, except so far as they have ceded it. \\
%            \pause
%            \vspace{30pt} \textbf{Dual Sovereignty}
%    \end{columns}
%\end{frame}
%
%\begin{frame}{John Taylor of Caroline, 1823}
%    \begin{columns}[c]
%        \column{0.5\textwidth}
%            \centering
%            \includegraphics[height=0.4\textheight]{img/john-taylor-of-caroline.png} \\
%        \column{0.5\textheight}
%            The \emph{consolidating school} contends that we have two
%            sovereignties; but that one is sovereign over the other; Mr.
%            Hamilton, that we have co-ordinate sovereignties, but that one is
%            made superlative.
%    \end{columns}
%\end{frame}

\begin{frame}{United States v. Gardner, 1997}
    \begin{columns}[c]
        \column{0.5\textwidth}
            \centering
            \includegraphics[width=\textwidth]{img/supcourt-building.jpg} \\
        \column{0.5\textheight}
            ``The state government and the federal government exercise \emph{concurrent jurisdiction} over the land.''
    \end{columns}
\end{frame}

\begin{frame}{Concurrent Jurisdiction Defined}
    \begin{columns}[c]
        \column{0.5\textwidth}
            \centering
            \includegraphics[height=0.4\textheight]{img/john-taylor-of-caroline.png} \\
        \column{0.5\textheight}
            \textbf{Concurrent Jurisdiction:} \\
            The \emph{consolidating school} contends that we have two
            sovereignties; but that one is sovereign over the other; Mr.
            Hamilton, that we have co-ordinate sovereignties, but that one is
            made superlative.
    \end{columns}
\end{frame}

\begin{frame}{Hamilton, 1788 newspaper}
    \begin{columns}[c]
        \column{0.5\textwidth}
            \centering
            \includegraphics[width=0.75\textwidth]{img/hamilton.png} \\
        \column{0.5\textheight}
            \only<1>{``Imperium in imperio is a political monster''}
            \only<2>{``[A sovereign within a sovereign] is a political monster''}
    \end{columns}
\end{frame}

%\begin{frame}{Kleppe v. New Mexico, 1976}
%    \begin{columns}[c]
%        \column{0.5\textwidth}
%            \centering
%            \includegraphics[width=0.75\textwidth]{img/octopus.png} \\
%            Department of the Interior \\
%        \column{0.5\textheight}
%            I have ``complete'' legislative power ``without limitation'' over Federal Lands!
%    \end{columns}
%\end{frame}
%
%\begin{frame}{FLPMA, 1976}
%    \begin{columns}[c]
%        \column{0.5\textwidth}
%            \centering
%            \includegraphics[width=0.75\textwidth]{img/octopus.png} \\
%            Department of the Interior \\
%        \column{0.5\textheight}
%            \emph{At \textbf{my} discretion} I will offer contracts to state
%            and local law enforcement officials to enforce Federal laws and
%            regulations on public lands.
%    \end{columns}
%\end{frame}
%
%\begin{frame}{Feds}
%    \begin{columns}[c]
%        \column{0.5\textwidth}
%            \centering
%            \includegraphics[width=0.75\textwidth]{img/fed-fish.png} \\
%        \column{0.5\textheight}
%            You states don't forget the ``Supremacy Clause''! We are more sovereign than you are!
%    \end{columns}
%\end{frame}
%
%\begin{frame}{Delegated powers, few and defined}
%    \begin{columns}[c]
%        \column{0.5\textwidth}
%            \centering
%            \includegraphics[width=0.75\textwidth]{img/constitution.png} \\
%        \column{0.5\textheight}
%            \begin{itemize}
%                \item Coin money
%                \item Declare war
%                \item Raise armies
%                \item Punish piracies
%                \item etc.
%            \end{itemize}
%    \end{columns}
%\end{frame}
%
%\begin{frame}
%    \begin{varblock}[0.8\textwidth]{Warning}
%        Federal Public Servants: Your powers are few and defined. If it is not on the list, you CANNOT do it!
%    \end{varblock}
%\end{frame}
%
%\begin{frame}
%    Delegated powers are NOT sovereignty! They are simply the powers the
%    sovereign states delegated (loaned) to their agent, the Federal Government.
%\end{frame}

\begin{frame}
    \begin{columns}[c]
        \column{0.5\textwidth}
            \centering
            \includegraphics[width=0.75\textwidth]{img/owl.png} \\
        \column{0.5\textheight}
            \Large{Do not confuse sovereignty with the powers which the sovereign delegates.}
    \end{columns}
\end{frame}

\begin{frame}{Jurisdiction Clause Family --- Guarantee}
    \centering
    \includegraphics[height=0.85\textheight]{img/family.png} \\
    \huge{\textbf{ \color{white}
        %\Put(25, 300){\begin{sideways}\colorbox{blue}{Admissions}\end{sideways}}
        %\Put(60, 300){\begin{sideways}\colorbox{blue}{Claims}\end{sideways}}
        %\Put(90, 300){\begin{sideways}\colorbox{blue}{Property}\end{sideways}}
        \Put(158, 300){\begin{sideways}\colorbox{blue}{Guarantee}\end{sideways}}
        %\Put(168, 300){\begin{sideways}\colorbox{blue}{Enclave}\end{sideways}}
        %\Put(200, 300){\begin{sideways}\colorbox{blue}{Engagements}\end{sideways}}
        %\Put(235, 300){\begin{sideways}\colorbox{blue}{Supremacy}\end{sideways}}
    }}
\end{frame}

\begin{frame}{Guarantee Clause}
    \centering
    \includegraphics[height=.7\textheight]{img/constitution.png} \\
    \textbf{Article IV, Section 4:} ``The United States shall \emph{guarantee} to every State in this Union a \emph{Republican} Form of Government''
\end{frame}

\begin{frame}{Principle of Constitutional Uniformity}
    \centering
    \includegraphics[height=.7\textheight]{img/us-map.png} \\
    \huge{Each state shall be a republic} \\
\end{frame}

\begin{frame}{Coyle v. Smith, 1911}
    \centering
    \includegraphics[width=0.75\textwidth]{img/sc-1905.png} \\
    ``\ldots the constitutional equality of the States is essential to the
    harmonious operation of the scheme upon which the Republic was organized.
    When that equality disappears we may remain free people, but the Union will
    not be the Union of the Constitution.''
\end{frame}

%\begin{frame}{What kind of republic?}
%    \begin{itemize}
%        \item Socialist Republic --- government control of all land
%        \pause
%        \item Aristocratic Republic --- wealthy individuals control government
%        \pause
%        \item Banana Republic --- deceitful and crafty individuals control government
%        \pause
%        \item ``the Republic''
%    \end{itemize}
%\end{frame}
%
%\begin{frame}
%    \centering
%    \large{``\ldots and to \textbf{the Republic} for which it stands\ldots''} \\
%    \includegraphics[width=.9\textwidth]{img/pledge.jpg} \\
%\end{frame}
%
%\begin{frame}{The Republic}
%    \begin{columns}[c]
%        \column{0.5\textwidth}
%            \begin{varblock}[0.8\textwidth]{Expert Witness}
%                \Large{In the \emph{compound republic} of America\ldots}
%            \end{varblock}
%        \column{0.5\textheight}
%            \centering
%            \includegraphics[width=0.9\textwidth]{img/madison.jpg} \\
%            James Madison, Father of the Constitution \\
%    \end{columns}
%\end{frame}
%
%\begin{frame}{The Republic of Republics}
%    \begin{columns}[c]
%        \column{0.5\textwidth}
%            \centering
%            \includegraphics[width=0.9\textwidth]{img/republic-of-republics.png} \\
%        \column{0.5\textheight}
%            \centering
%            Bernard Janin Sage \\
%            1878, 606 pages \\
%    \end{columns}
%\end{frame}
%
%\begin{frame}{Who Owns the West?}
%    \centering
%    \includegraphics[width=.7\textwidth]{img/federal-lands.png} \\
%    \pause
%    \Large{The Constitutional requirement to \textbf{dispose} of public land was abandoned!} \\
%\end{frame}

\begin{frame}
    \begin{block}{In Our Republican Form of Government}
        At the time of statehood, the federal government \emph{relinquished jurisdiction} and proprietary interest over the land with two exceptions:
        \begin{itemize}
            \item Indian lands
            \item Enclaves
        \end{itemize}
        Admission of new states is ``on an equal footing with the original states in \textbf{all respects whatever}.'' \\
    \end{block}
\end{frame}

\begin{frame}{Equal Footing Doctrine}
    \centering
    \includegraphics[height=0.95\textheight]{img/federal-lands.png} \\
\end{frame}

\begin{frame}{Guarantee Clause has been severely violated}
    \centering
    \includegraphics[height=.7\textheight]{img/constitution.png} \\
    \textbf{Article IV, Section 4:} ``The United States shall \emph{guarantee} to every State in this Union a \emph{Republican} Form of Government''
\end{frame}

\begin{frame}{REMEMBER! Johnson v. McIntosh, 1823}
    \centering
    \includegraphics[width=0.75\textwidth]{img/sc-1905.png} \\
    ``And the \emph{vacant soil is to be disposed of} by that organ of the government which has the constitutional power to dispose of the national domain.'' \\
\end{frame}

\begin{frame}{Law School}
    \centering
    \includegraphics[height=.7\textheight]{img/law-school.png} \\
    \pause
    \large{``Don't read the Constitution. It is too confusing. Here we study case law.''}
\end{frame}

\begin{frame}{United States v. Gratiot, 1840}
    \centering
    \includegraphics[width=0.65\textwidth]{img/sc-1905.png} \\
%    \only<1>{
%    ``\ldots and this power [over the territorial lands] is vested in Congress
%    without limitation, and has been considered the \emph{foundation} upon which the
%    \emph{territorial} government rests.''
%    }
%    \only<2>{
    ``\ldots and this power [over the territorial lands] is vested in Congress
    without limitation {\huge\color{red} \textbf{,}} and has been considered the \emph{foundation} upon which the
    territorial government rests.''
%    }
\end{frame}

\begin{frame}{Kleppe v. New Mexico, 1976}
    \centering
    \includegraphics[width=0.65\textwidth]{img/sc-1976.png} \\
    ``We have repeatedly observed that the power over the public lands is thus entrusted to Congress without limitation{ \huge \color{red} \textbf{.}}''
\end{frame}

\begin{frame}{Constitutional Principles of Jurisdiction}
    \begin{columns}[onlytextwidth]
        \column{0.5\textwidth}
            { \large
            Gratiot, 1840: Congress has unlimited power over \textbf{\emph{pre-statehood territorial lands}}. \\
            \vspace{20pt}
            Kleppe, 1976: Congress has unlimited power over \textbf{\emph{``public lands'' within a state}}. \\
            }
        \column{0.5\textwidth}
            \centering
            \includegraphics[width=0.9\textwidth]{img/red-light.png} \\
    \end{columns}
\end{frame}

\begin{frame}{Jurisdiction Clause Family --- Property}
    \centering
    \includegraphics[height=0.85\textheight]{img/family.png} \\
    \huge{\textbf{ \color{white}
        %\Put(25, 300){\begin{sideways}\colorbox{blue}{Admissions}\end{sideways}}
        %\Put(60, 300){\begin{sideways}\colorbox{blue}{Claims}\end{sideways}}
        \Put(105, 300){\begin{sideways}\colorbox{blue}{Property}\end{sideways}}
        %\Put(134, 300){\begin{sideways}\colorbox{blue}{Guarantee}\end{sideways}}
        %\Put(168, 300){\begin{sideways}\colorbox{blue}{Enclave}\end{sideways}}
        %\Put(200, 300){\begin{sideways}\colorbox{blue}{Engagements}\end{sideways}}
        %\Put(235, 300){\begin{sideways}\colorbox{blue}{Supremacy}\end{sideways}}
    }}
\end{frame}

\begin{frame}{August 30, 1787}
    \centering
    \includegraphics[width=0.85\textwidth]{img/convention.png} \\
\end{frame}

\begin{frame}{The ``Trust Compacts'' or Engagements, August 30, 1787}
    \begin{columns}[onlytextwidth]
        \column{0.5\textwidth}
            \vspace{10pt}
            \begin{itemize}
                \item The ``Resolution of 1780'' formed the basis upon which Congress was \textbf{required to dispose} of territorial and public lands.
                \item Art. of Confederation, Article IX: ``no State shall be deprived of territory for the benefit of the United States.''
                \item Other state land cession compacts
            \end{itemize}
        \column{0.5\textwidth}
            \centering
            \includegraphics[width=0.85\textwidth]{img/convention.png} \\
    \end{columns}
\end{frame}

\begin{frame}{The ``Trust Compacts'' or Engagements, August 30, 1787}
    \begin{columns}[onlytextwidth]
        \column{0.5\textwidth}
            { \Large It would be best to leave everything on that subject in ``statu quo.'' }
        \column{0.5\textwidth}
            \centering
            \includegraphics[width=0.85\textwidth]{img/convention.png} \\
    \end{columns}
\end{frame}

\begin{frame}{The ``Trust Compacts'' or Engagements, August 30, 1787}
    \begin{columns}[onlytextwidth]
        \column{0.5\textwidth}
            \centering
            \includegraphics[height=0.4\textheight]{img/morris-portrait.png} \\
            { \tiny  Gouvernor Morris } \\
            { \Large ``I make a motion'' \ldots for a property clause} \\
        \column{0.5\textwidth}
            \centering
            \includegraphics[width=0.85\textwidth]{img/convention.png} \\
    \end{columns}
\end{frame}

\begin{frame}{Property Clause}
   \centering
   \includegraphics[height=.7\textheight]{img/constitution.png} \\
   \textbf{Article IV, sec. 3, cl. 2:} ``Congress shall have Power to \emph{dispose of} and make all needful Rules and Regulations respecting the Territory or other \emph{Property} belonging to the United States;''
\end{frame}

\begin{frame}{Property Clause, in \emph{untampered}, plain English}
   \centering
   \includegraphics[height=.7\textheight]{img/constitution.png} \\
   When a territory becomes a state, Congress shall have power to make all needful rules and regulations to \emph{dispose} of former territorial lands and other property in an orderly manner.
\end{frame}

\begin{frame}{Benner v. Porter, 1850}
    \centering
    \includegraphics[width=0.75\textwidth]{img/sc-1905.png} \\
    ``\ldots we are satisfied that \ldots the act of Congress admitting the Territory of Florida, as a State, into the Union \ldots displaced the territorial government, and \emph{abrogated all its powers and jurisdiction.}''
\end{frame}

\begin{frame}{Benner v. Porter, 1850}
    \centering
    \includegraphics[width=0.75\textwidth]{img/sc-1905.png} \\
    ``\ldots the territorial government was displaced, abrogated, every part of it; and that \emph{no power of jurisdiction} existed within her limits, except that derived from State authority.''
\end{frame}

\begin{frame}{Jurisdiction Clause Family --- Enclave}
    \centering
    \includegraphics[height=0.85\textheight]{img/family.png} \\
    \huge{\textbf{ \color{white}
        %\Put(25, 300){\begin{sideways}\colorbox{blue}{Admissions}\end{sideways}}
        %\Put(60, 300){\begin{sideways}\colorbox{blue}{Claims}\end{sideways}}
        %\Put(105, 300){\begin{sideways}\colorbox{blue}{Property}\end{sideways}}
        %\Put(134, 300){\begin{sideways}\colorbox{blue}{Guarantee}\end{sideways}}
        \Put(198, 300){\begin{sideways}\colorbox{blue}{Enclave}\end{sideways}}
        %\Put(200, 300){\begin{sideways}\colorbox{blue}{Engagements}\end{sideways}}
        %\Put(235, 300){\begin{sideways}\colorbox{blue}{Supremacy}\end{sideways}}
    }}
\end{frame}

\begin{frame}{Enclave Clause}
   \centering
   \includegraphics[height=.6\textheight]{img/constitution.png} \\
   \textbf{Article I, sec. 8, cl. 17:} ``Congress shall have power \ldots To
   exercise exclusive Legislation \ldots over all Places \emph{purchased by the
   Consent of the Legislature of the State} in which the Same shall be, for the
   Erection of Forts, Magazines, Arsenals, dock-Yards and other needful
   Buildings.''
\end{frame}

\begin{frame}{Who Owns the West?}
    \centering
    \includegraphics[width=.7\textwidth]{img/federal-lands.png} \\
    \Large{Were the ``western lands'' ``purchased by the Consent of the Legislature of the State''?} \\
\end{frame}

\begin{frame}{Fundamental Principle}
   \centering
   \includegraphics[height=.6\textheight]{img/constitution.png} \\
   \textbf{Fundamental Principle:} By the Constitution only the States can
   create parks, forest reserves, wildlife refuges, monuments, etc. The
   federal government may hold \textbf{no} land in a state for
   non-constitutional purposes.
\end{frame}

\begin{frame}{Federal Attorneys, 2012}
    \begin{columns}[c]
        \column{0.5\textwidth}
            \centering
            Incorrect\ldots \\
            incorrect\ldots \\
            misguided\ldots \\
            Property Clause\ldots \\
            Property Clause\ldots \\
            Property Clause\ldots \\
            Supremacy Clause\ldots \\
            numerous court decisions\ldots \\
            federal laws\ldots \\
            NO STATE CONSENT NECESSARY! \\
        \column{0.5\textwidth}
            \centering
            \includegraphics[height=0.5\textheight]{img/elmer-dickens.png} \\
            Amanda Marshall \\
            Kenneth Paur \\
    \end{columns}
\end{frame}

\begin{frame}{Federal Attorneys, 2012}
\begin{table}[h]
\centering
\begin{tabular}{p{.20\textwidth} p{.70\textwidth}} 
    { \tiny Attorney Dickens } &
        \multirow{2}{.7\textwidth}{ ``There are many US Supreme Court cases going
        back \emph{over a hundred years} where the court has found that Congress has
        broad authority under the Property Clause.''} \\
        { \includegraphics[height=0.15\textheight]{img/elmer-dickens.png} } & \\
    { \tiny Attorney Marshall } & 
        \multirow{2}{.7\textwidth}{``Supreme Court precedent for the \emph{past one hundred years} is clear: under the Property Clause Congress may delegate rule making authority to the federal agencies.'' } \\
        { \includegraphics[height=0.15\textheight]{img/amanda-marshall.png} } & \\
    { \tiny Attorney Paur } & 
        \multirow{2}{.7\textwidth}{``Authority of the Secretary of Agriculture\ldots recognized for \emph{over one hundred years}\ldots Property Clause\ldots Property Clause\ldots'' (20X)} \\
        { \includegraphics[height=0.15\textheight]{img/elmer-dickens.png} } & \\
\end{tabular}
\end{table}
\end{frame}

%\begin{frame}{Constitutional Fortress}
%    \centering
%    \includegraphics[width=0.75\textwidth]{img/fortress.png} \\
%    \pause
%    \Put(10,180){\includegraphics[width=0.3\textwidth]{img/family.png}}
%    \pause
%    { \color{red} \huge{Nationalist Assault} \\}
%    \pause
%    \Put(40,180){\includegraphics[width=0.3\textwidth]{img/property-ram.png}}
%    \pause
%    \Put(140,180){\includegraphics[width=0.3\textwidth]{img/supremacy-ram.png}}
%\end{frame}
%
%\begin{frame}{Constitutional Fortress}
%    \centering
%    \includegraphics[width=0.75\textwidth]{img/fortress.png} \\
%    \Put(10,180){\includegraphics[width=0.3\textwidth]{img/family.png}}
%    { \color{red} \huge{Nationalist Assault} \\}
%    \Put(90,220){\includegraphics[width=0.3\textwidth]{img/blam.png}}
%    \Put(40,180){\includegraphics[width=0.3\textwidth]{img/property-ram.png}}
%    \Put(140,180){\includegraphics[width=0.3\textwidth]{img/supremacy-ram.png}}
%\end{frame}
%
%\begin{frame}{The Federal Position}
%    \begin{columns}[c]
%        \column{0.3\textwidth}
%            \centering
%            \includegraphics[width=0.75\textwidth]{img/elmer-dickens.png} \\
%            Attorney Kenneth Paur \\
%            Nov. 10, 2011 \\
%            Property Clause x 20 \\
%        \column{0.7\textwidth}
%            The Property Clause provides separate and independent authority for
%            Congress to make rules respecting federal property, without
%            qualification or limitation. \ldots  Under the Supremacy Clause,
%            federal statutes enacted pursuant to the Property clause preempt
%            conflicting state laws.
%    \end{columns}
%\end{frame}

\begin{frame}
    \begin{varblock}[.8\textwidth]{Problem}
        The wellspring of USFS and BLM power is judicial misconstruction of the Property Clause.
    \end{varblock}
\end{frame}

\begin{frame}
    \begin{columns}[c]
        \column{0.5\textwidth}
            \centering
            \includegraphics[height=0.75\textheight]{img/htwwl.png} \\
            2000, 446 pages \\
        \column{0.5\textwidth}
            \centering
            \includegraphics[height=0.35\textheight]{img/william-hayward.png} \\
            { \large The Theft and Usurpation of State's Property Rights }
    \end{columns}
\end{frame}

\begin{frame}
    \begin{columns}[c]
        \column{0.5\textwidth}
            \centering
            \includegraphics[height=0.35\textheight]{img/william-hayward.png} \\
            William Hayward, Page 384 \\
        \column{0.5\textwidth}
            ``In this case [Utah Power \& Light v. United States, 1917], \emph{as in all others I have examined}, [federal land ownership] was not decided; it was \emph{assumed}.''
    \end{columns}
\end{frame}

\begin{frame}
    \begin{columns}[c]
        \column{0.5\textwidth}
            \centering
            \includegraphics[height=0.75\textheight]{img/htwwl.png} \\
            2000 \\
        \column{0.5\textwidth}
            Page 149 \\
            {\Large Eight Supreme Court cases support: \emph{Federal jurisdiction can only occur \textbf{{\color{red} with the consent}} of the legislature of the state. }}
    \end{columns}
\end{frame}

\begin{frame}{Orwellian Doublethink}
    \centering
    \includegraphics[width=.85\textwidth]{img/orwell.png} \\
    \Put(5,85){\includegraphics[width=.3\textwidth]{img/1984.png}}
\end{frame}

\begin{frame}{Orwellian Doublethink}
    \vspace{1pt}
    { \Large
    \textbf{Enclave Clause:} State consent \emph{is} necessary for the Federal Government to build a Post Office within a state \\
    \vspace{20pt}
    \pause
    \textbf{Contemporary Interpretation:} \emph{No} state consent is necessary for the Federal Government to establish a 1.6 million acre national monument within a State. \\
    }
    \pause
    \Put(100,200){\includegraphics[with=.4\textwidth]{img/red-light.png}}
\end{frame}

%\begin{frame}
%    \begin{varblock}[.9\textwidth]{Fundamental Rule of Constitutional Interpretation}
%        An enumerated delegation of limited power will NOT be followed by a general grant of power.
%    \end{varblock}
%\end{frame}

\begin{frame}{Jurisdiction Clause Family --- Property}
    \centering
    \includegraphics[height=0.85\textheight]{img/family.png} \\
    \huge{\textbf{ \color{white}
        %\Put(25, 300){\begin{sideways}\colorbox{blue}{Admissions}\end{sideways}}
        %\Put(60, 300){\begin{sideways}\colorbox{blue}{Claims}\end{sideways}}
        \Put(108, 300){\begin{sideways}\colorbox{blue}{Property}\end{sideways}}
        %\Put(134, 300){\begin{sideways}\colorbox{blue}{Guarantee}\end{sideways}}
        %\Put(168, 300){\begin{sideways}\colorbox{blue}{Enclave}\end{sideways}}
        %\Put(200, 300){\begin{sideways}\colorbox{blue}{Engagements}\end{sideways}}
        %\Put(235, 300){\begin{sideways}\colorbox{blue}{Supremacy}\end{sideways}}
    }}
\end{frame}

\begin{frame}{Property Clause}
   \centering
   \includegraphics[height=.7\textheight]{img/constitution.png} \\
   \textbf{Article IV, sec. 3, cl. 2:} ``Congress shall have Power to \emph{dispose of} and make all needful Rules and Regulations respecting the Territory or other \emph{Property} belonging to the United States;''
\end{frame}

\begin{frame}{\st{Property Clause} Disposal Clause}
   \centering
   \includegraphics[height=.7\textheight]{img/constitution.png} \\
   \textbf{Article IV, sec. 3, cl. 2:} ``Congress shall have Power to \emph{dispose of} and make all needful Rules and Regulations respecting the Territory or other \emph{Property} belonging to the United States;''
\end{frame}

\begin{frame}{``Disposal'' Clause, in \emph{untampered}, plain English}
   \centering
   \includegraphics[height=.7\textheight]{img/constitution.png} \\
   \only<1>{When a territory becomes a state, Congress shall have power to make all needful rules and regulations to \emph{dispose} of former territorial lands and other property in an orderly manner. \\ }
   \only<2>{This gives Congress the power to \emph{dispose} of federal lands, not the power \emph{not to dispose} of them. \\}
\end{frame}
%
%\section{End of section 2}
%\frame{\mysectionpage}

%\begin{frame}{News Release}
%    \begin{block}{March 23, 2012}
%    SALT LAKE CITY - Today, Utah Governor Gary R. Herbert signed House Bill
%    148, which demands the federal government make good on the promises made in
%    the 1894 Enabling Act to extinguish title to federal lands in Utah. 
%    \end{block}
%\end{frame}
%
%\begin{frame}
%    \begin{columns}[c]
%        \column{0.5\textwidth}
%            \centering
%            \includegraphics[width=0.75\textwidth]{img/thomas-paine.png} \\
%            Thomas Paine, 1737 - 1809 \\
%        \column{0.5\textwidth}
%            ``All power exercised over a nation must have some beginning. It
%            must either be delegated or assumed. There are no other sources.
%            All delegated power is trust, and all assumed power is usurpation.
%            Time does not alter the nature and quality of either.''
%    \end{columns}
%\end{frame}
%
%\begin{frame}
%   \centering
%   \includegraphics[height=.99\textheight]{img/the-end.png} \\
%\end{frame}

%\section{Life after Kleppe}
%\frame{\mysectionpage}

\begin{frame}{Gouvernor Morris and the Property Clause}
    \centering
    \includegraphics[height=.8\textheight]{img/morris-full.png} \\
\end{frame}

\begin{frame}
    \begin{columns}[onlytextwidth]
        \column{0.5\textwidth}
            \textbf{Letter, December 4, 1803} \\
            Morris, by his own admission, had attempted to concoct a deception,
            if not a fraud, hidden in duplicitous wording of the Property
            Clause.

        \column{0.5\textwidth}
            \centering
            \includegraphics[width=0.95\textwidth]{img/morris-portrait.png} \\
            Gouvernor Morris \\
    \end{columns}
\end{frame}

\begin{frame}
    \begin{columns}[onlytextwidth]
        \column{0.5\textwidth}
            \textbf{Letter, December 4, 1803} \\
            ``I went as far as circumstances would permit\ldots Had it been
            more pointedly expressed, a strong opposition would have been made.''

        \column{0.5\textwidth}
            \centering
            \includegraphics[width=0.95\textwidth]{img/morris-portrait.png} \\
            Gouvernor Morris \\
    \end{columns}
\end{frame}

\begin{frame}
    \begin{columns}[onlytextwidth]
        \column{0.5\textwidth}
            \centering
            \includegraphics[width=0.95\textwidth]{img/james-kent.png} \\
            James Kent, 1826 \\
            \emph{Commentaries on American Law}, 4 vols. \\
            \small{\emph{as quoted in Statehood, p. 115}} \\

        \column{0.5\textwidth}
            { \large ``[Judges] roam at large in the trackless field of their own imaginations.''}
    \end{columns}
\end{frame}

%\begin{frame}{1976, Property Clause Revisited}
%    \centering
%    \includegraphics[width=0.75\textwidth]{img/sc-1976.png} \\
%\end{frame}
%
%\begin{frame}{1976, Property Clause Revisited}
%    \centering
%    \includegraphics[width=0.75\textwidth]{img/sc-1976.png} \\
%    ``\ldots make all needful rules and regulations\ldots'' \\
%    \Put(260,60){\includegraphics[width=.2\textwidth]{img/sherlock.png}}
%\end{frame}
%
%\begin{frame}{``Emanations from Penumbrae''}
%   \centering
%   \includegraphics[height=.7\textheight]{img/constitution.png} \\
%   \Put(200,080){\includegraphics[width=.4\textwidth]{img/eyes.png}}
%\end{frame}
%
%\begin{frame}{``Emanations from Penumbrae''}
%   \centering
%   \includegraphics[height=.7\textheight]{img/fire.png} \\
%   \Put(200,080){\includegraphics[width=.4\textwidth]{img/eyes.png}}
%\end{frame}
%
%\begin{frame}{``Emanations from Penumbrae''}
%    \begin{varblock}[.9\textwidth]{}
%        \textbf{penumbrae:} noun, The partially shaded outer region of the shadow cast by an opaque object.
%    \end{varblock}
%\end{frame}

\begin{frame}{Kleppe v. New Mexico, 1976}
    \centering
    \includegraphics[width=0.75\textwidth]{img/sc-1976.png} \\
    { \large Federal legislative power over public lands is ``complete'' and
    ``without limitation.'' \\ }
%    \Put(260,160){\includegraphics[width=.2\textwidth]{img/sherlock.png}}
\end{frame}

%\begin{frame}
%    \begin{columns}[onlytextwidth]
%        \column{0.5\textwidth}
%            { \large Imagine Gouvernor Morris now\ldots }
%        \column{0.5\textwidth}
%            \centering
%            \includegraphics[width=0.95\textwidth]{img/morris-smile.png} \\
%    \end{columns}
%\end{frame}

\begin{frame}{Jurisdiction Clause Family --- Supremacy}
    \centering
    \includegraphics[height=0.85\textheight]{img/family.png} \\
    \huge{\textbf{ \color{white}
        %\Put(25, 300){\begin{sideways}\colorbox{blue}{Admissions}\end{sideways}}
        %\Put(60, 300){\begin{sideways}\colorbox{blue}{Claims}\end{sideways}}
        %\Put(90, 300){\begin{sideways}\colorbox{blue}{Property}\end{sideways}}
        %\Put(134, 300){\begin{sideways}\colorbox{blue}{Guarantee}\end{sideways}}
        %\Put(168, 300){\begin{sideways}\colorbox{blue}{Enclave}\end{sideways}}
        %\Put(200, 300){\begin{sideways}\colorbox{blue}{Engagements}\end{sideways}}
        \Put(235, 300){\begin{sideways}\colorbox{blue}{Supremacy}\end{sideways}}
    }}
\end{frame}

\begin{frame}{Supremacy Clause}
   \centering
   \includegraphics[height=.7\textheight]{img/constitution.png} \\
   \textbf{Article VI, clause 2:} ``This Constitution, and the Laws of the
   United States \emph{which shall be made in Pursuance thereof}\ldots shall be
   the supreme Law of the Land.''
\end{frame}

%\begin{frame}{United States Senate and House of Representatives}
%    \centering
%    \includegraphics[width=0.75\textwidth]{img/house-senate-fail.png} \\
%\end{frame}

\begin{frame}{Kleppe v. New Mexico, 1976}
    \centering
    \includegraphics[width=0.65\textwidth]{img/sc-1976.png} \\
    Congress may enact legislation governing federal lands pursuant to the
    property clause and ``when Congress so acts, federal legislation
    necessarily overrides conflicting state laws under the \emph{supremacy
    clause}.''
\end{frame}

\begin{frame}{Kohl v. United States, 1875}
    \centering
    \includegraphics[width=0.75\textwidth]{img/sc-1905.png} \\
    { \large ``The consent of a State can never be a condition precedent to [the exercise of federal powers].'' }
\end{frame}

\begin{frame}{Ft. Leavenworth R.R. v. Lowe, 1885}
    \centering
    \includegraphics[width=0.75\textwidth]{img/sc-1905.png} \\
    The Kohl decision is ``good law,'' ``a doctrine authoritatively declared.'' Congress may exercise jurisdiction over the places it might acquire \emph{without State consent} provided that those places are ``used as a means to carry out the purposes of the government.''
\end{frame}

%\begin{frame}{Continuum of \textbf{\emph{Sliding}} Judicial Opinions}
%    \Put(20,20){\Large \begin{turn}{50} Johnson v. McIntosh, 1823 \end{turn}}
%    \Put(40,10){\Large \begin{turn}{50} United States v. Gratiot, 1840 \end{turn}}
%    \Put(50,10){\Large \begin{turn}{50} Benner v. Porter, 1850 \end{turn}}
%    \Put(60,10){\Large \begin{turn}{50} Friedman v. Goodwin, 1856 \end{turn}}
%    \Put(80,30){\Large \begin{turn}{50} Kohl v. United States, 1875 \end{turn}}
%    \Put(110,45){\Large \begin{turn}{50} Ft. Leavenworth v. Lowe, 1885 \end{turn}}
%    \Put(120,45){\Large \begin{turn}{50} South Carolina v. United States, 1905 \end{turn}}
%    \Put(130,40){\Large \begin{turn}{50} Collins v. Yosemite Park, 1938 \end{turn}}
%    \Put(170,30){\Large \begin{turn}{50} Kleppe v. New Mexico, 1976 \end{turn}}
%    \Put(190,-115){\Large \begin{turn}{50} FPLMA, 1976 \end{turn}}
%    {
%        \color{red}
%        \thicklines
%        \Put(-20, -130){ \line(1, 0){250} }
%    }
%%    \pause
%%    \Put(-30,100){\includegraphics[width=.3\textwidth]{img/hh-coin.png}}
%%    \pause
%%    \Put(260,0){\includegraphics[height=.4\textheight]{img/fasces.png}}
%%    \pause
%%    \Put(30,-240){\includegraphics[width=.5\textwidth]{img/coffin-clause.png}}
%\end{frame}
{%
\setlength{\fboxsep}{0pt}%
\setlength{\fboxrule}{1pt}%
\begin{frame}[t]{Contiuum of Sliding Judicial Opinions}
    \centering
%    Johnson v. McIntosh, 1823
    \only<1>{ \fbox{\includegraphics[height=0.6\textheight]{img/sliding/1.png}} \\ 
        ``And the vacant soil is to be disposed of by that organ of the
        government which has the constitutional power to dispose of the
        national domain.'' \\
    }
%    United States v. Gratiot, 1840
    \only<2>{ \fbox{\includegraphics[height=0.6\textheight]{img/sliding/2.png}} \\ 
        ``\ldots and this power [over the territorial lands] is vested in
        Congress without limitation, and has been considered the foundation
        upon which the territorial government rests.'' \\
    }
%    Benner v. Porter, 1850
    \only<3>{ \fbox{\includegraphics[height=0.6\textheight]{img/sliding/3.png}} \\ 
        ``\ldots we are satisfied that \ldots the act of Congress admitting the
        Territory of Florida, as a State, into the Union \ldots displaced the
        territorial government, and \emph{abrogated all its powers and
        jurisdiction.}'' \\
    }
%    Friedman v. Goodwin, 1856
    \only<4>{ \fbox{\includegraphics[height=0.6\textheight]{img/sliding/4.png}} \\ 
        ``On [California's] admission into the Union, she became the owner of
        all the public land not disposed of by law of congress.'' \\
    }
%    Kohl v. United States, 1875
    \only<5>{ \fbox{\includegraphics[height=0.6\textheight]{img/sliding/5.png}} \\ 
         ``The consent of a State can never be a condition precedent to [the
         exercise of federal powers].'' \\
    }
%    Ft. Leavenworth v. Lowe, 1885
    \only<6>{ \fbox{\includegraphics[height=0.6\textheight]{img/sliding/6.png}} \\ 
        The Kohl decision is ``good law,'' ``a doctrine authoritatively
        declared.'' Congress may exercise jurisdiction over the places it might
        acquire without State consent provided that those places are
        ``used as a means to carry out the purposes of the government.'' \\
    }
%    South Carolina v. United States, 1905
    \only<7>{ \fbox{\includegraphics[height=0.6\textheight]{img/sliding/7.png}} \\ 
        ``The Constitution is a written instrument. As such, its meaning does
        not alter. That which it meant when it was adopted, it means now.'' \\
    }
%    Collins v. Yosemite Park, 1938
    \only<8>{ \fbox{\includegraphics[height=0.6\textheight]{img/sliding/8.png}} \\ 
        ``The national government may, `by virtue of its sovereignty' acquire
        lands within the borders of states by eminent domain and without their
        consent.'' (referring to Kohl) \\
    }
%    Kleppe v. New Mexico, 1976
    \only<9>{ \fbox{\includegraphics[height=0.6\textheight]{img/sliding/9.png}} \\ 
        ``We have repeatedly observed that the power over the public lands is
        thus entrusted to Congress without limitation.'' \\
    }
\end{frame}
}%


%\begin{frame}{Here lies the Clause Family}
%    \centering
%    \includegraphics[width=.6\textwidth]{img/coffin.png} \\
%    { \large
%    Born March 4, 1789 \\
%    Died by Strangulation, 1976 \\
%    }
%\end{frame}

\begin{frame}{126 Days after Kleppe: FLPMA}
    \centering
    \includegraphics[width=.8\textwidth]{img/congress.png} \\
    Federal Lands Policy and Management Act \\
    { \large ``\ldots that it is the policy of the United States that the
    public lands be retained in Federal ownership\ldots'' }
\end{frame}

\begin{frame}
    \includegraphics[width=.3\textwidth]{img/stop.png}
    \hspace{10pt}
    \includegraphics[width=.3\textwidth]{img/think.png}
    \hspace{10pt}
    \includegraphics[width=.3\textwidth]{img/wrong-way.png} \\
    \vspace{20pt}
    \begin{block}{Consequences of Kleppe}
        Public lands within the States have been rendered jurisdictionally
        equivalent to pre-statehood Territories, the District of Columbia, and
        federal enclaves!
    \end{block}
\end{frame}

\begin{frame}{Jurisdiction Clause Family --- Claims}
    \centering
    \includegraphics[height=0.85\textheight]{img/family.png} \\
    \huge{\textbf{ \color{white}
        %\Put(25, 300){\begin{sideways}\colorbox{blue}{Admissions}\end{sideways}}
        \Put(70, 300){\begin{sideways}\colorbox{blue}{Claims}\end{sideways}}
        %\Put(90, 300){\begin{sideways}\colorbox{blue}{Property}\end{sideways}}
        %\Put(134, 300){\begin{sideways}\colorbox{blue}{Guarantee}\end{sideways}}
        %\Put(168, 300){\begin{sideways}\colorbox{blue}{Enclave}\end{sideways}}
        %\Put(200, 300){\begin{sideways}\colorbox{blue}{Engagements}\end{sideways}}
        %\Put(235, 300){\begin{sideways}\colorbox{blue}{Supremacy}\end{sideways}}
    }}
\end{frame}

\begin{frame}{Claims Clause}
   \centering
   \includegraphics[height=.7\textheight]{img/constitution.png} \\
   \only<1>{\textbf{Article IV, sec. 3, cl. 2:} ``\ldots nothing in this Constitution shall be so construed as to Prejudice any \emph{Claims}\ldots of any particular State.''}
   \only<2>{\textbf{Article IV, sec. 3, cl. 2:} \st{``\ldots nothing in this Constitution shall be so construed as to Prejudice any \emph{Claims}\ldots of any particular State.''}}
\end{frame}

\begin{frame}{Engagements Clause}
   \centering
   \includegraphics[height=.7\textheight]{img/constitution.png} \\
   \only<1>{\textbf{Article VI, clause 1:} ``All\ldots \emph{Engagements} entered into,
   before the Adoption of this Constitution, shall be as valid against the
   United States under this Constitution, as under the Confederation.''}
   \only<2>{\textbf{Article VI, clause 1:} \st{``All\ldots \emph{Engagements} entered into,
   before the Adoption of this Constitution, shall be as valid against the
   United States under this Constitution, as under the Confederation.''}}
\end{frame}

\begin{frame}{Engagements Clause}
   \centering
   \includegraphics[height=.7\textheight]{img/constitution.png} \\
   By an act of Congress, the Constitution was amended by a change in policy to
   abandon the requirement to dispose of territorial and public land.
\end{frame}

\begin{frame}{Jurisdiction Clause Family}
    \centering
    \includegraphics[height=0.85\textheight]{img/family.png} \\
    \huge{\textbf{ \color{white}
        \Put(25, 300){\begin{sideways}\colorbox{blue}{Admissions}\end{sideways}}
        \Put(60, 300){\begin{sideways}\colorbox{blue}{Claims}\end{sideways}}
        \Put(90, 300){\begin{sideways}\colorbox{blue}{Property}\end{sideways}}
        \Put(134, 300){\begin{sideways}\colorbox{blue}{Guarantee}\end{sideways}}
        \Put(168, 300){\begin{sideways}\colorbox{blue}{Enclave}\end{sideways}}
        \Put(200, 300){\begin{sideways}\colorbox{blue}{Engagements}\end{sideways}}
        \Put(235, 300){\begin{sideways}\colorbox{blue}{Supremacy}\end{sideways}}
    }}
\end{frame}

\begin{frame}{Jurisdiction Clause Family}
    \centering
    \includegraphics[height=0.85\textheight]{img/family.png} \\
    \huge{\textbf{ \color{white}
        \Put(17, 300){ \begin{sideways}\colorbox{red}{\st{Admissions}}\end{sideways}}
        \Put(52, 300){ \begin{sideways}\colorbox{red}{\st{Claims}}\end{sideways}}
        \Put(82, 300){ \begin{sideways}\colorbox{red}{\st{Property}}\end{sideways}}
        \Put(134, 300){\begin{sideways}\colorbox{red}{\st{Guarantee}}\end{sideways}}
        \Put(168, 300){\begin{sideways}\colorbox{red}{\st{Enclave}}\end{sideways}}
        \Put(200, 300){\begin{sideways}\colorbox{red}{\st{Engagements}}\end{sideways}}
        \Put(235, 300){\begin{sideways}\colorbox{red}{\st{Supremacy}}\end{sideways}}
    }}
\end{frame}

\begin{frame}{REMEMBER!}
    \includegraphics[height=0.75\textheight]{img/colonies.png} \\
    \Put(200,300){\includegraphics[width=0.4\textwidth]{img/constitution.png}}
    \Put(195,420){\textbf{List of delegated powers}}
    \Put(200,60){\includegraphics[width=0.3\textwidth]{img/national-seal.png}}
    \Put(225,170){\textbf{\Large{Agent}}}
\end{frame}

\begin{frame}
    \begin{columns}[c]
        \column{0.5\textwidth}
            \Large{What if the agent does something that is not on the list?}
        \column{0.5\textheight}
            \centering
            \includegraphics[width=0.95\textwidth]{img/kid-question.png} \\
    \end{columns}
\end{frame}

\begin{frame}{Resolutions of 1798}
    \begin{columns}[c]
        \column{0.5\textwidth}
            \Large{Each state is duty bound to ``interpose'' between the people and the federal government.}
        \column{0.5\textheight}
            \centering
            \includegraphics[height=0.4\textheight]{img/madison.jpg} \\
            { \tiny James Madison \\ }
            \includegraphics[height=0.4\textheight]{img/jefferson.png} \\
            { \tiny Thomas Jefferson \\ }
    \end{columns}
\end{frame}

\begin{frame}{James Madison, Virginia Resolution, 1798}
    That this Assembly doth explicitly and peremptorily declare, that it views
    the powers of the federal government as resulting from the \textbf{
    \color{blue}compact} to which the states are parties, as limited by
    the plain sense and intention of the instrument constituting that
    \textbf{\color{blue}compact}, as no further valid than they are
    authorized by the grants enumerated in that\textbf{ \color{blue}
    compact}; and that, in case of a deliberate, palpable, and dangerous
    exercise of other powers, not granted by the said\textbf{ \color{blue}
    compact}, the states, who are parties thereto, have the right, and are
    in duty bound, to\textbf{ \color{red}interpose}, for arresting the
    progress of the evil, and for maintaining, within their respective
    limits, the authorities, rights and liberties, appertaining to them.  
\end{frame}

\begin{frame}{Interposition}
    \only<1>{\vspace{10pt}}
    \only<2>{Sovereign states and subdivisions are duty-bound to draw the line}
    \includegraphics[height=0.75\textheight]{img/colonies.png} \\
    \Put(200,60){\includegraphics[width=0.3\textwidth]{img/national-seal.png}}
    \Put(100,200){\includegraphics[height=0.8\textheight]{img/line.png}}
\end{frame}

\begin{frame}
    \begin{block}{IMPORTANT!}
        { \huge Interposition is \emph{not} violent. It is \emph{peaceful refusal to allow an agent to exceed its authority}.}
    \end{block}
\end{frame}

\begin{frame}{Example \#1 --- Sheriff Phil Barney}
    \begin{columns}[onlytextwidth]
        \column{0.5\textwidth}
            November 7, 2000 \\
            \vspace{10pt}
            { \Large ``Ranchers and sheriff defy feds, take back seized herd of cattle'' }

        \column{0.5\textwidth}
            \centering
            \includegraphics[width=0.75\textwidth]{img/phil-barney.png} \\
            Sheriff Phil Barney, Sevier County, Utah \\
    \end{columns}
\end{frame}

\begin{frame}{Example \#1 --- Sheriff Phil Barney}
    \Put(0,90){\includegraphics[height=0.5\textheight,keepaspectratio=true,width=0.4\textwidth]{img/barney-ranchers.png}}
    \Put(210,-280){\includegraphics[height=0.5\textheight,keepaspectratio=true,width=0.4\textwidth]{img/blm.png}}
    \Put(100,-180){\includegraphics[height=0.5\textheight,keepaspectratio=true,width=0.4\textwidth]{img/phil-barney.png}}
\end{frame}

\begin{frame}{Example \#2 --- Sheriff Tony de Meo}
    \Put(0,90){\includegraphics[height=0.5\textheight,keepaspectratio=true,width=0.4\textwidth]{img/de-meo-ranch.png}}
    \Put(0,-220){\includegraphics[height=0.5\textheight,keepaspectratio=true,width=0.4\textwidth]{img/de-meo-person.png}}
    \Put(210,-280){\includegraphics[height=0.5\textheight,keepaspectratio=true,width=0.4\textwidth]{img/blm.png}}
    \Put(100,-180){\includegraphics[height=0.5\textheight,keepaspectratio=true,width=0.4\textwidth]{img/tony-de-meo.png}}
\end{frame}

\begin{frame}
    \begin{columns}[onlytextwidth]
        \column{0.5\textwidth}
            \centering
            \includegraphics[width=0.95\textwidth]{img/carrie-dann.png} \\
            Carrie Dann, Nevada Rancher \\

        \column{0.5\textwidth}
            ``I always tell people I was free until 1973, before the BLM came.''
    \end{columns}
\end{frame}

\begin{frame}{Example \#3 --- Sheriff Benny House}
    \begin{columns}[onlytextwidth]
        \column{0.5\textwidth}
            \centering
            \includegraphics[width=0.75\textwidth]{img/benny-house.png} \\
            Sheriff Benny House \\
            Otero County, New Mexico \\

        \column{0.5\textwidth}
            September 17, 2011 \\
            \vspace{10pt}
            { \Large Tree Party Rebellion }
    \end{columns}
\end{frame}

\begin{frame}{Example \#3 --- Sheriff Benny House}
    \begin{columns}[onlytextwidth]
        \column{0.5\textwidth}
            \centering
            \includegraphics[width=0.75\textwidth]{img/bill-cox.png} \\
            Bill Cox, Utah Rancher \\

        \column{0.5\textwidth}
            ``When I was a boy the United States Forest Service was our best friend. Now the Forest Service and BLM are our worst enemies!''
    \end{columns}
\end{frame}

\begin{frame}{Example \#4 --- Sheriff Brad Rogers}
    \begin{columns}[onlytextwidth]
        \column{0.5\textwidth}
            \centering
            \includegraphics[width=0.75\textwidth]{img/brad-rogers.png} \\
            Sheriff Brad Rogers \\
            Elkhart County, Indiana \\

        \column{0.5\textwidth}
            December 17, 2011 \\
            \vspace{10pt}
            { \Large Sheriff Brad Rogers of Indiana Stands Up Against the FDA  \\}
            \vspace{10pt}
            Forest Grove Dairy the focus of federal investigation by the FDA
    \end{columns}
\end{frame}

\begin{frame}{Example \#4 --- Sheriff Brad Rogers}
    \Put(0,90){\includegraphics[height=0.5\textheight,keepaspectratio=true,width=0.4\textwidth]{img/hochstetler-dairy.png}}
    \Put(0,-20){ David Hochstetler }
    \Put(210,-280){\includegraphics[height=0.5\textheight,keepaspectratio=true,width=0.4\textwidth]{img/fda.png}}
    \Put(100,-180){\includegraphics[height=0.5\textheight,keepaspectratio=true,width=0.4\textwidth]{img/brad-rogers.png}}
\end{frame}

\begin{frame}{Utah --- Bound by Oath}
    \begin{columns}[c]
        \column{0.5\textwidth}
            \centering
            \includegraphics[height=0.9\textheight]{img/utah-capitol.png} \\
        \column{0.5\textheight}
            \centering
            \includegraphics[height=0.9\textheight]{img/utah-sheriffs.png} \\
    \end{columns}
\end{frame}

% XXX Fix hokey pictures here
\begin{frame}
    \begin{columns}[c]
        \column{0.5\textwidth}
            \only<1>{The \textbf{federal} courts have held that the power to declare \textbf{federal} laws unconstitutional lies with the \textbf{federal} courts, not with the states.}
            \only<2>{\centering \includegraphics[width=0.8\textwidth]{img/fox-chickens.png} \\}
        \column{0.5\textheight}
            \centering
            \includegraphics[width=0.9\textwidth]{img/judge-clip.png} \\
    \end{columns}
\end{frame}

\begin{frame}{The Supreme Court}
    \centering
    \includegraphics[width=0.9\textwidth]{img/supreme.png} \\
    \only<1>{\large{We are the living Constitution: What we say is what the Constitution means. \\}}
    \only<2>{\large{Interposition is an illegal defiance of Constitutional authority.\\}}
\end{frame}

\begin{frame}{Correction, Please}
    \centering
    \includegraphics[width=0.9\textwidth]{img/capitol.png} \\
    \large{\emph{\textbf{Interposition}} is a peaceful defiance of illegal unconstitutional authority. \\}
\end{frame}

\begin{frame}
    \centering
    \large{When the legislative, executive, and judicial branches have betrayed us, the states must \textbf{draw the line}, as sentinels of the liberties of the people. \\}
    \includegraphics[height=0.7\textheight]{img/kids.png} \\
\end{frame}

\begin{frame}{Massachussetts, 1809}
    \begin{columns}[c]
        \column{0.5\textwidth}
            \centering
            \includegraphics[width=0.9\textwidth]{img/massachussetts.png} \\
            Massachussetts Legislature, 1809
        \column{0.5\textheight}
            The Embargo Act is ``unconstitutional'' and ``not legally binding on the citizens of this state''. \\
            \vspace{20pt}
            \pause
            \Large{\color{red}Nullification}
    \end{columns}
\end{frame}

\begin{frame}{``Nullification'', 2010}
    \begin{columns}[c]
        \column{0.5\textwidth}
            \centering
            \includegraphics[width=0.9\textwidth]{img/nullification.png} \\
        \column{0.5\textheight}
            \centering
            \only<1>{\includegraphics[width=0.9\textwidth]{img/thomas-woods.png} \\}
            \only<2>{\includegraphics[height=0.4\textheight]{img/ed-phillips.png} \\ Sheriff Ed Phillips \\ Millard County, Utah \\
            \vspace{10pt}
                \large{``I wish to hell that I had had this book in my hands 20 years ago when we started the Western Sheriffs Association.''}
            }
    \end{columns}
\end{frame}

%\begin{frame}
%    \centering
%    \includegraphics[height=.9\textheight]{img/territorial-imperative} \\
%    Bill Howell, Bill Redd --- 2005, 524 pages \\
%\end{frame}
%
%
%\begin{frame}{Evolutionary Socialism, 1910}
%    \centering
%    \includegraphics[height=.9\textheight]{img/remold-it-small.png} \\
%\end{frame}

\begin{frame}{The World's Plan for Happiness}
    \centering
    \includegraphics[height=.9\textheight]{img/remold-it.jpg} \\
\end{frame}

\begin{frame}{Fundamental Principles of Socialism}
    \centering
    \includegraphics[height=.9\textheight]{img/fabian.png} \\
\end{frame}

\begin{frame}{Fundamental Principles of Socialism}
    \centering
    1. Government ownership \emph{or control} of all land \\
    \includegraphics[height=.9\textheight]{img/federal-lands.png} \\
\end{frame}

\begin{frame}{Government Ownership of Nevada}
    \centering
    \includegraphics[width=.9\textwidth]{img/nevada.jpg} \\
\end{frame}

\begin{frame}{Fundamental Principles of Socialism}
    \centering
    2. Government ownership \emph{or control} of major industries\\
    \includegraphics[width=.9\textwidth]{img/industries.jpg} \\
\end{frame}

\begin{frame}{Fundamental Principles of Socialism}
    \centering
    3. Government control over labor \\
    \includegraphics[width=.9\textwidth]{img/labor.jpg} \\
\end{frame}

\begin{frame}{Fundamental Principles of Socialism}
    \centering
    4. Government ownership \emph{or control} of communications and transportation \\
    \includegraphics[width=.9\textwidth]{img/transportation.jpg} \\
\end{frame}

\begin{frame}{Fundamental Principles of Socialism}
    \centering
    5. Government control of all credit \\
    \includegraphics[width=.9\textwidth]{img/credit.png} \\
\end{frame}

\begin{frame}{Fundamental Principles of Socialism}
    \centering
    6. Government control of all insurance \\
    \includegraphics[width=.9\textwidth]{img/insurance.jpg} \\
\end{frame}

\begin{frame}{Fundamental Principles of Socialism}
    \centering
    7. Government control of the educational system \\
    \includegraphics[height=.9\textheight]{img/education.jpg} \\
\end{frame}

\begin{frame}{Fundamental Principles of Socialism}
    \centering
    8. Elimination of the significance of the family \\
    \includegraphics[height=.9\textheight]{img/gay.jpg} \\
\end{frame}

\begin{frame}{Fundamental Principles of Socialism}
    \centering
    9. Elimination of the significance of religion \\
    \only<1>{\includegraphics[width=.9\textwidth]{img/10-commandments.jpg} \\}
    \only<2>{\includegraphics[width=.9\textwidth]{img/religion.png} \\}
\end{frame}

\begin{frame}{Fundamental Principles of Socialism}
    \centering
    10. Establishment of a minimum wage \\
    \includegraphics[width=.9\textwidth]{img/wage.jpg} \\
\end{frame}

\begin{frame}{Fundamental Principles of Socialism}
    \centering
    11. A universal system of pensions \\
    \includegraphics[width=.9\textwidth]{img/pensions.jpg} \\
\end{frame}

\begin{frame}{Fundamental Principles of Socialism}
    \centering
    12. Justified use of force if necessary to attain socialistic goals
    \includegraphics[width=.9\textwidth]{img/justified-force.jpg} \\
\end{frame}

\begin{frame}{Fundamental Principles of Socialism}
    \centering
    13. Graduated income tax \\
    \includegraphics[width=.9\textwidth]{img/graduated-tax.jpg} \\
    \only<2>{``Tell us where you are, and we'll send you money!''}
\end{frame}

\begin{frame}{Fundamental Principles of Socialism}
    \begin{enumerate}
        \item Government ownership or control of all land
        \item Government ownership or control of major industries
        \item Government control over labor
        \item Government ownership or control of communications and transportation
        \item Government control of all credit
        \item Government control of all insurance
        \item Government control of the educational system
        \item Elimination of the significance of the family
        \item Elimination of the significance of religion
        \item Establishment of a minimum wage
        \item A universal system of pensions
        \item Justified use of force where necessary to attain socialistic goals
        \item Graduated income tax
    \end{enumerate}
\end{frame}

%\begin{frame}{1910}
%    \begin{columns}[onlytextwidth]
%        \column{0.5\textwidth}
%            \centering
%            \includegraphics[width=0.95\textwidth]{img/remold-hammers.png} \\
%            \only<1>{ \color{white} \huge ``Gradual, silent encroachment'' }
%            \only<2>{ \huge ``Gradual, silent encroachment'' }
%
%        \column{0.5\textwidth}
%            \centering
%            \includegraphics[width=0.95\textwidth]{img/wolf.png} \\
%    \end{columns}
%\end{frame}

\begin{frame}{Gradual Encroachment}
    \centering
    \includegraphics[width=0.65\textwidth]{img/french-pres.png} \\
    Val\'{e}ry Giscard d'Estaing \\
\end{frame}

\begin{frame}{Sustainable Development}
    \begin{columns}[onlytextwidth]
        \column{0.5\textwidth}
            \centering
            \includegraphics[width=.75\textwidth]{img/agenda-21.png} \\

        \column{0.5\textwidth}
            \centering
            \includegraphics[width=.75\textwidth]{img/ghwbush.png} \\
            President George H. W. Bush, 1992 \\
    \end{columns}
\end{frame}

\begin{frame}{U.N. Biodiversity Treaty}
    \centering
    \includegraphics[width=0.85\textwidth]{img/biodiversity.png} \\
    Never ratified by the U.S. Senate \\
\end{frame}

\begin{frame}{Wildlands Project}
    \centering
    \includegraphics[width=0.75\textwidth]{img/wildlands.png} \\
    The Wildlands Project is designed to remove humans from rural lands and allow forests to grow out of control, even when forest fires threaten human life and property. \\
\end{frame}

\begin{frame}{President's Council on Sustainable Development}
    \begin{columns}[onlytextwidth]
        \column{0.5\textwidth}
            \centering
            \includegraphics[width=.75\textwidth]{img/clinton.png} \\
            President Bill Clinton \\

        \column{0.5\textwidth}
            ``We will implement Agenda 21 by administrative laws and executive orders.''
    \end{columns}
\end{frame}

\begin{frame}{Antiquities Act, 1906}
    \begin{columns}[onlytextwidth]
        \column{0.5\textwidth}
            \centering
            Congress has not acted quickly enough to protect federal lands. \\
            \vspace{20pt}
            { \Large{``Monument Power'' 22x}}

        \column{0.5\textwidth}
            \centering
            \includegraphics[width=.75\textwidth]{img/clinton.png} \\
            President Bill Clinton \\

    \end{columns}
\end{frame}

\begin{frame}
    \centering
    \includegraphics[width=0.75\textwidth]{img/salazar.png} \\
    January 10, 2010: Dept. of Interior Secretary Ken Salazar announces
    \textbf{new mining near the Grand Canyon will be banned on a million acres} \\
\end{frame}

\begin{frame}{National Monuments}
    \centering
    \includegraphics[width=0.75\textwidth]{img/monuments.png} \\
    \begin{itemize}
        \item Prohibit natural resource development: timber, mining, oil, natural gas, etc.
        \item Close roads
        \item Establish wilderness areas
        \item Remove dams
    \end{itemize}
\end{frame}

% XXX Dean wilson  is the first on the left side. Greg Hagwood is in the middle, right next to some lady. Indicate that somehow on their quotes.
\begin{frame}{Yreka, California --- October 22, 2011}
    \centering
    \includegraphics[width=0.95\textwidth]{img/yreka.png} \\
    \vspace{20pt}
    \only<1>{
    The Federal Government pays environmental groups to sue the Federal
    Government to stop the use of your property! For example, \$36 million was
    paid to 19 environmental groups in 9 years in 19 states. \\
	Karen-Budd Falen, Property Rights Lawyer \\ }
    \only<2>{ ``I had spent a good part of my life enforcing the penal code,
    but not understanding my oath of office.'' \\ Sheriff Dean Wilson \\ }
    \only<3>{ ``A giant has been awakened, and [the federal bureaucracy] didn't count on that.'' \\
    Sheriff Greg Hagwood \\ }
\end{frame}

\begin{frame}{Sheriff Glenn Palmer}
    \begin{columns}[onlytextwidth]
        \column{0.5\textwidth}
            \centering
            \includegraphics[width=0.75\textwidth]{img/glenn-palmer-letter.png}
            \\ { \tiny March 31, 2011 }
        \column{0.5\textwidth}
            \centering
            \includegraphics[width=0.75\textwidth]{img/glenn-palmer.png}
            \\ Sheriff Glenn Palmer
            \\ Grant County, Oregon
    \end{columns}
\end{frame}

\begin{frame}{Sheriff Glenn Palmer}
    \begin{columns}[onlytextwidth]
        \column{0.5\textwidth}
            \centering
            Where does the United States Forest Service get its Constitutional authority to have law enforcement officers in Grant County?
        \column{0.5\textwidth}
            \centering
            \includegraphics[width=0.75\textwidth]{img/glenn-palmer.png}
            \\ Sheriff Glenn Palmer
            \\ Grant County, Oregon
    \end{columns}
\end{frame}

\begin{frame}{Sheriff Glenn Palmer}
    \begin{columns}[onlytextwidth]
        \column{0.5\textwidth}
            \centering
            \includegraphics[width=0.75\textwidth]{img/glenn-palmer.png}
            \\ Sheriff Glenn Palmer
            \\ Grant County, Oregon

        \column{0.5\textwidth}
            ``One response that I have received in writing is that their
            authority is given through the Cooperative Policing Agreement that
            this agency has signed in the past. Upon asking for clarification
            and a second request, the response was that I needed to check with
            my District Attorney. Neither response in my opinion is adequate.''
    \end{columns}
\end{frame}

\begin{frame}{USDA, December 20, 2010}
    \begin{columns}[onlytextwidth]
        \column{0.5\textwidth}
            \centering
            \includegraphics[width=0.75\textwidth]{img/usda-letter.png} \\

        \column{0.5\textwidth}
            5301 -- Forest Service \textbf{law enforcement authority} \\
            5301.1 -- Laws
            \indent 1. The Act of June 4, 1897 \\
            \indent 8. The Act of 2004 \\
            \indent 1. Title 7 U.S.C \\
    \end{columns}
\end{frame}

\begin{frame}{Sheriff Glenn Palmer}
    \begin{columns}[onlytextwidth]
        \column{0.5\textwidth}
            \centering
            \includegraphics[width=0.75\textwidth]{img/glenn-palmer.png}
            \\ Sheriff Glenn Palmer
            \\ Grant County, Oregon

        \column{0.5\textwidth}
            { \Large
                ``Constitutional authority,'' \textbf{NOT} ``Law enforcement authority''
                \pause
                Shouldn't their ``constitutional authority'' be referenced to the Constitution?
            }
    \end{columns}
\end{frame}

%\begin{frame}
%    \centering
%    { \Large
%        LEGAL UPDATE FOR LAW ENFORCEMENT \\
%        NEW MEXICO SHERIFFS' \& POLICE ASSOCIATION \\
%        NOVEMBER 10, 2011 \\
%        ********* \\
%        FOREST SERVICE LAW ENFORCEMENT AUTHORITIES \\
%        \vspace{20pt}
%        Prepared by:
%        Kenneth D. Paur, Assistant Regional Attorney \\
%        U.S.D.A. Office of the General Counsel \\
%        Ogden, Utah \\
%        (801) 625-5440 kenneth.paur@oge.usda.gov \\
%        (Twenty pages) \\
%    }
%\end{frame}

\begin{frame}{Legal Update for Law Enforcement, New Mexico}
    \begin{varblock}[.9\textwidth]{}
        ``Property Clause''\ldots
        ``Property Clause''\ldots
        ``Property Clause''\ldots
        Supreme Court rejected the need for State consent\ldots
        Supreme Court rejected\ldots
        ``Supremacy Clause''\ldots
    \end{varblock}
\end{frame}

\begin{frame}{Legal Update for Law Enforcement, New Mexico}
    \begin{varblock}[.9\textwidth]{Page 3}
        ``The federal enabling statute imposed several conditions for
        recognition of a state and admission into the Union, including the
        requirement that \textbf{\emph{the state and its citizens forever
        disclaim any right or title to the federal lands within the state}}.''
    \end{varblock}
\end{frame}

\begin{frame}{Legal Update for Law Enforcement, New Mexico}
    \begin{varblock}[.9\textwidth]{Page 20}
        ``Tens of thousands of criminal cases have been investigated and
        prosecuted, and convictions obtained for violation of laws relating to
        the national forests during this time [from 1891 to 2011].''
    \end{varblock}
\end{frame}

\begin{frame}{Sheriff Gil Gilbertson}
    \begin{columns}[onlytextwidth]
        \column{0.5\textwidth}
            \centering
            Where does the United States Forest Service's authority come from?
            \\ { \tiny October, 2011 }
        \column{0.5\textwidth}
            \centering
            \includegraphics[width=0.75\textwidth]{img/gil-gilbertson.png}
            \\ Sheriff Gil Gilbertson
            \\ Josephine County, Oregon
    \end{columns}
\end{frame}

\begin{frame}{James Madison}
    \begin{columns}[onlytextwidth]
        \column{0.5\textwidth}
            \centering
            \includegraphics[width=0.75\textwidth]{img/madison.jpg} \\
        \column{0.5\textwidth}
            ``Since the general civilization of mankind, I believe there are
            more instances of the abridgment of the freedom of the people by
            \emph{gradual and silent encroachments} of those in power, than by
            violent and sudden usurpation.''
    \end{columns}
\end{frame}

\begin{frame}{Sheriff Gil Gilbertson}
    \begin{columns}[onlytextwidth]
        \column{0.5\textwidth}
            \centering
            Where does the United States Forest Service's authority come from? \\
            \vspace{20pt}
            \textbf{ \Large \color{red}
                From ``gradual and silent encroachments''
            }
        \column{0.5\textwidth}
            \centering
            \includegraphics[width=0.75\textwidth]{img/gil-gilbertson.png}
            \\ Sheriff Gil Gilbertson
            \\ Josephine County, Oregon
    \end{columns}
\end{frame}

\begin{frame}{The Truth Is:}
    \large
    The U.S. Forest Service and BLM have NO \textbf{Constitutional authority}.
    By judicial misconstruction of the Property clause, social do-gooders have
    reinvented the structure of the United States' government.
\end{frame}

\begin{frame}{The Salt Lake Tribune}
    \begin{block}{March 8, 2011}
    A bill that would ensure Utah students learn the U.S. is a compound constitutional republic --- not a democracy --- has passed both
    Houses of the Legislature and is now headed to the governor for his signature.
    \end{block}
\end{frame}
%
%\begin{frame}{Utah Governor's Race, 2012}
%    \centering
%    \includegraphics[width=0.8\textwidth,height=.5\textwidth,keepaspectratio=true]{img/gov.png} \\
%    { \Large These three candidates for Utah governor said they would \emph{actively fight to regain state control of federal land} in Utah, refuse federal funding, and apply ``nullification.'' (Jan, 2012) \\ }
%\end{frame}
% XXX Include stuff about Utah's enabling act somewhere. Also pencil drawings, but make the Octopus have "State" on each arm

%\begin{frame}{What is this?}
%    \centering
%    \includegraphics[height=0.95\textheight]{img/rack.png} \\
%\end{frame}
%
%\begin{frame}
%    \begin{minipage}{0.2\textwidth}
%        \centering
%        \only<1,2>{\includegraphics[width=0.9\textwidth]{img/white-wolf.png} \\}
%        \only<4->{\includegraphics[width=0.9\textwidth]{img/wolf.png} \\}
%        \only<1-3>{\includegraphics[width=0.9\textwidth]{img/white-wolf.png} \\}
%        \only<1-4>{\includegraphics[width=0.9\textwidth]{img/white-turtle.png} \\}
%        \only<5->{\includegraphics[width=0.9\textwidth]{img/turtle.png} \\}
%    \end{minipage}
%    \begin{minipage}{0.5\textwidth}
%        \centering
%        \includegraphics[height=0.9\textheight]{img/rack.png}
%        \only<6->{\Put(-110,210){\includegraphics[width=.6\textwidth]{img/us-map-stretched.png}}}
%    \end{minipage}
%    \begin{minipage}{0.2\textwidth}
%        \centering
%        \only<2->{\includegraphics[width=0.9\textwidth]{img/national-seal.png} \\ }
%        \only<1,2->{\includegraphics[height=0.5\textheight]{img/white-fasces.png}}
%        \only<3->{\includegraphics[height=0.5\textheight]{img/fasces.png}}
%    \end{minipage}
%\end{frame}
%
%\begin{frame}
%    \begin{minipage}{0.2\textwidth}
%        \centering
%        \includegraphics[width=0.9\textwidth]{img/catch.png} \\
%        \only<4->{\Put(-35,0){ \colorbox{white}{\color{red} \Large \textbf{Judicial}}}}
%    \end{minipage}
%    \begin{minipage}{0.5\textwidth}
%        \centering
%        \includegraphics[height=0.9\textheight]{img/rack.png}
%        \Put(-110,210){\includegraphics[width=.6\textwidth]{img/us-map-stretched.png}}
%        \only<2->{\Put(-120,120){ \colorbox{white}{\color{red} \Large \textbf{Legislative}}}}
%        \only<3->{\Put(-90,290){ \colorbox{white}{\color{red} \Large \textbf{Executive}}}}
%    \end{minipage}
%    \begin{minipage}{0.2\textwidth}
%        \centering
%        \includegraphics[width=0.9\textwidth]{img/white-seal.png} \\ 
%    \end{minipage}
%\end{frame}
%
%\begin{frame}{``Ratchet Effect''}
%    \begin{columns}[onlytextwidth]
%        \column{0.5\textwidth}
%            \centering
%            \includegraphics[width=0.9\textwidth]{img/catch.png} \\
%
%        \column{0.5\textwidth}
%            \begin{itemize}
%                \item Unconstitutional laws
%                \item USC
%                \item Regulations
%                \item Executive orders
%                \item Court decisions
%                \item \ldots and so on
%            \end{itemize}
%            
%    \end{columns}
%\end{frame}

\begin{frame}{Federal Register Growth}
    \centering
    \includegraphics[width=0.75\textwidth]{img/fed-reg-graph.png} \\
    { \large 200 pages per day \\}
\end{frame}

\begin{frame}{Federal Register Growth}
    The 2009 edition of the Code of Federal Regulations contained 163,333 pages in 226 individual books
\end{frame}

\begin{frame}{Federal Register Growth}
    \centering
    \includegraphics[width=0.75\textwidth]{img/fed-reg-bookshelf.png} \\
    { \large ``A \emph{few} copies of the Code of Federal Regulations'' \\ }
\end{frame}

\begin{frame}{James Madison}
    \begin{columns}[onlytextwidth]
        \column{0.5\textwidth}
            \centering
            \includegraphics[width=0.75\textwidth]{img/madison.jpg} \\
        \column{0.5\textwidth}
            ``It will be of little avail to the people that the laws are made
            by men of their own choice if the laws be so voluminous that they
            cannot be read, or so incoherent that they cannot be understood.''
    \end{columns}
\end{frame}


%\begin{frame}{Constitutional Sheriffs, 2012}
%\begin{table}[h]
%\centering
%\begin{tabular}{cccc} 
%    \includegraphics[width=0.75\textwidth,height=.24\textheight,keepaspectratio=true]{img/s1.png} &
%    \includegraphics[width=0.75\textwidth,height=.24\textheight,keepaspectratio=true]{img/s2.png} &
%    \includegraphics[width=0.75\textwidth,height=.24\textheight,keepaspectratio=true]{img/s3.png} &
%    \includegraphics[width=0.75\textwidth,height=.24\textheight,keepaspectratio=true]{img/s4.png} \\
%    \includegraphics[width=0.75\textwidth,height=.24\textheight,keepaspectratio=true]{img/tony-de-meo.png} &
%    \includegraphics[width=0.75\textwidth,height=.24\textheight,keepaspectratio=true]{img/s6.png} &
%    \includegraphics[width=0.75\textwidth,height=.24\textheight,keepaspectratio=true]{img/s7.png} &
%    \includegraphics[width=0.75\textwidth,height=.24\textheight,keepaspectratio=true]{img/s8.png} \\
%    &
%    \includegraphics[width=0.75\textwidth,height=.24\textheight,keepaspectratio=true]{img/s9.png} &
%    \includegraphics[width=0.75\textwidth,height=.24\textheight,keepaspectratio=true]{img/gil-gilbertson.png} \\
%\end{tabular}
%\end{table}
%\end{frame}

%

\begin{frame}{James Madison}
    \begin{columns}[onlytextwidth]
        \column{0.5\textwidth}
            \centering
            \includegraphics[width=0.75\textwidth]{img/madison-2.png} \\
            James Madison, Federalist 46 \\
        \column{0.5\textwidth}
            The State governments are to protect the Union of States from ``encroachments of the Federal Government.''
    \end{columns}
\end{frame}

\begin{frame}
    \centering
    \includegraphics[height=0.95\textheight]{img/fed-papers-copy.jpg} \\
\end{frame}

\begin{frame}{Nationalists}
\begin{table}[h]
\centering
\begin{tabular}{lcccccc} 
    \includegraphics[width=0.75\textwidth,height=.3\textheight,keepaspectratio=true]{img/hamilton-portrait.png} &
    \includegraphics[width=0.75\textwidth,height=.3\textheight,keepaspectratio=true]{img/lincoln-portrait.png} &
    \includegraphics[width=0.75\textwidth,height=.3\textheight,keepaspectratio=true]{img/morris-portrait.png} \\
    Alexander Hamilton & 
    Abraham Lincoln &
    Gouvernor Morris \\
    \includegraphics[width=0.75\textwidth,height=.3\textheight,keepaspectratio=true]{img/story-portrait.png} &
    \includegraphics[width=0.75\textwidth,height=.3\textheight,keepaspectratio=true]{img/marshall-portrait.png} &
    \includegraphics[width=0.75\textwidth,height=.3\textheight,keepaspectratio=true]{img/daniel-webster-portrait.png} \\
    Joseph Story &
    John Marshall &
    Daniel Webster &
\end{tabular}
\end{table}
\end{frame}

\begin{frame}{States' Rights Advocates}
\begin{table}[h]
\centering
\begin{tabular}{lcccccc} 
    \includegraphics[width=0.75\textwidth,height=.3\textheight,keepaspectratio=true]{img/jefferson.png} &
    \includegraphics[width=0.75\textwidth,height=.3\textheight,keepaspectratio=true]{img/luther-martin.png} &
    \includegraphics[width=0.75\textwidth,height=.3\textheight,keepaspectratio=true]{img/patrick-henry.png} \\
    Thomas Jefferson &
    Luther Martin &
    Patrick Henry \\
    \includegraphics[width=0.75\textwidth,height=.3\textheight,keepaspectratio=true]{img/john-taylor-of-caroline.png} &
    \includegraphics[width=0.75\textwidth,height=.3\textheight,keepaspectratio=true]{img/st-george-tucker.png} &
    \includegraphics[width=0.75\textwidth,height=.3\textheight,keepaspectratio=true]{img/george-mason.png} \\
    John Taylor of Caroline &
    St. George Tucker &
    George Mason \\
\end{tabular}
\end{table}
\end{frame}

\begin{frame}{Interposition}
    \only<1>{\vspace{10pt}}
    \only<2>{Sovereign states and subdivisions are duty-bound to draw the line}
    \includegraphics[height=0.75\textheight]{img/colonies.png} \\
    \Put(200,60){\includegraphics[width=0.3\textwidth]{img/national-seal.png}}
    \Put(120,220){\includegraphics[height=0.8\textheight]{img/line.png}}
\end{frame}

\begin{frame}{What does it mean to be bound by oath to support \emph{this} constitution?}
    \begin{columns}[onlytextwidth]
        \column{0.5\textwidth}
            \centering
            \includegraphics[width=0.75\textwidth]{img/bork-book.jpg} \\
            1990 \\

        \column{0.5\textwidth}
            ``\emph{This} Constitution, not one they make up, is to bind [public servants].''

    \end{columns}
\end{frame}

\begin{frame}{What does it mean to be bound by oath to support \emph{this} constitution?}
    \begin{columns}[onlytextwidth]
        \column{0.5\textwidth}
            \centering
            \includegraphics[width=0.75\textwidth]{img/bork-book2.png} \\
            2003 \\

        \column{0.5\textwidth}
            {\large
                ``The Constitution is simply gone.'' \\
                \vspace{20pt}
                The government is not following the Constitution.
            }

    \end{columns}
\end{frame}

\begin{frame}{Where did the Constitution go?}
    \begin{columns}[onlytextwidth]
        \column{0.5\textwidth}
            \centering
            \includegraphics[width=0.75\textwidth]{img/joe-sobran.png} \\
            Joe Sobran \\

        \column{0.5\textwidth}
            \textbf{How Tyranny Came to America} \\
            \vspace{20pt}
            {\large
                ``The United States Constitution poses no serious threat to our form of government today.''
            }

    \end{columns}
\end{frame}

%\begin{frame}{Constitutional Sheriffs, 2012}
%\begin{table}[h]
%\centering
%\begin{tabular}{cccc} 
%    \includegraphics[width=0.75\textwidth,height=.24\textheight,keepaspectratio=true]{img/s1.png} &
%    \includegraphics[width=0.75\textwidth,height=.24\textheight,keepaspectratio=true]{img/s2.png} &
%    \includegraphics[width=0.75\textwidth,height=.24\textheight,keepaspectratio=true]{img/s3.png} &
%    \includegraphics[width=0.75\textwidth,height=.24\textheight,keepaspectratio=true]{img/s4.png} \\
%    \includegraphics[width=0.75\textwidth,height=.24\textheight,keepaspectratio=true]{img/tony-de-meo.png} &
%    \includegraphics[width=0.75\textwidth,height=.24\textheight,keepaspectratio=true]{img/s6.png} &
%    \includegraphics[width=0.75\textwidth,height=.24\textheight,keepaspectratio=true]{img/s7.png} &
%    \includegraphics[width=0.75\textwidth,height=.24\textheight,keepaspectratio=true]{img/s8.png} \\
%    &
%    \includegraphics[width=0.75\textwidth,height=.24\textheight,keepaspectratio=true]{img/s9.png} &
%    \includegraphics[width=0.75\textwidth,height=.24\textheight,keepaspectratio=true]{img/gil-gilbertson.png} \\
%\end{tabular}
%\end{table}
%\end{frame}

\begin{frame}{Bound by Oath}
    \centering
    \includegraphics[width=0.75\textwidth]{img/sheriff-oath.png} \\
\end{frame}
%
%
%\begin{frame}{Bound by Oath}
%    \centering
%    \includegraphics[width=0.75\textwidth]{img/sheriff-oath.png} \\
%\end{frame}

\begin{frame}
   \centering
   \includegraphics[height=.99\textheight]{img/the-end.png} \\
\end{frame}

\begin{frame}
   \centering
   \includegraphics[height=.99\textheight]{img/end1.png} \\
\end{frame}

\begin{frame}
   \centering
   \includegraphics[height=.99\textheight]{img/end2.png} \\
\end{frame}

\begin{frame}
   \centering
   \includegraphics[height=.99\textheight]{img/end3.png} \\
\end{frame}

\begin{frame}
   \centering
   \includegraphics[height=.99\textheight]{img/end4.png} \\
\end{frame}

%\begin{frame}{Verse 1}
% \huge{
%    \centering
%    Amazing Grace, how sweet the sound, \\
%    that saved a wretch like me. \\
%    I once was lost but now am found, \\
%    was blind, but now, I see. \\
%}
%\end{frame}
%
%\begin{frame}{Verse 2}
% \huge{
%    \centering
%    T'was Grace that taught my heart to fear, \\
%    And Grace, my fears relieved. \\
%    How precious did that Grace appear \\
%    the hour I first believed. \\
%}
%\end{frame}
%
%\begin{frame}{Verse 3}
% \huge{
%    \centering
%    Through many dangers, toils and snares \\
%    I have already come. \\
%    T'was Grace that brought us safe thus far \\
%    and Grace will lead me home. \\
%}
%\end{frame}
%
%\begin{frame}{Verse 4}
% \huge{
%    \centering
%    The Lord has promised good to me. \\
%    His word my hope secures. \\
%    He will my shield and portion be \\
%    as long as life endures. \\
%}
%\end{frame}
%
%\begin{frame}{Verse 5}
% \huge{
%    \centering
%    Yea, when this flesh and heart shall fail \\
%    and mortal life shall cease, \\
%    I shall possess within the veil \\
%    a life of joy and peace. \\
%}
%\end{frame}
%
%\begin{frame}{Verse 6}
% \huge{
%    \centering
%    When we've been here ten thousand years \\
%    bright shining as the sun, \\
%    We've no less days to sing God's praise \\
%    than when we've first begun. \\
%}
%\end{frame}
%
%\begin{frame}{Verse 7}
% \huge{
%    \centering
%    Amazing Grace, how sweet the sound, \\
%    that saved a wretch like me. \\
%    I once was lost but now am found, \\
%    was blind, but now, I see. \\
%}
%\end{frame}

\end{document}
